%%%%%%%% ICML 2025 EXAMPLE LATEX SUBMISSION FILE %%%%%%%%%%%%%%%%%

\documentclass[sigconf,anonymous,review]{article}
\AtBeginDocument{%
	\providecommand\BibTeX{{%
			Bib\TeX}}}

% Recommended, but optional, packages for figures and better typesetting:
\usepackage{microtype}
\usepackage{graphicx}
\usepackage{subfigure}
\usepackage{booktabs} % for professional tables
\usepackage{dsfont}

% hyperref makes hyperlinks in the resulting PDF.
% If your build breaks (sometimes temporarily if a hyperlink spans a page)
% please comment out the following usepackage line and replace
% \usepackage{icml2025} with \usepackage[nohyperref]{icml2025} above.
\usepackage{hyperref}
%\usepackage{icml2025}

% Attempt to make hyperref and algorithmic work together better:
\newcommand{\theHalgorithm}{\arabic{algorithm}}

% Use the following line for the initial blind version submitted for review:
% \usepackage{icml2025}

% If accepted, instead use the following line for the camera-ready submission:
%\usepackage[accepted]{icml2025}

% For theorems and such
\usepackage{amsmath}
\usepackage{amssymb}
%\usepackage{mathtools}
%\usepackage{amsthm}

% if you use cleveref..
\usepackage[capitalize,noabbrev]{cleveref}

%\newcommand{\eat}[1]{}


%%%%commm

% Todonotes is useful during development; simply uncomment the next line
%    and comment out the line below the next line to turn off comments
%\usepackage[disable,textsize=tiny]{todonotes}
%\usepackage[textsize=tiny]{todonotes}


\usepackage{latexsym}
\usepackage{amsfonts}
\usepackage{amsmath}
%\usepackage{amssymb}
\usepackage{color}
\usepackage{epsfig}
\usepackage{xspace}
\usepackage{graphicx}
%\usepackage{subcaption}
\usepackage{cleveref}
%\usepackage{times}
\usepackage{booktabs}
\usepackage{diagbox}
\usepackage{multirow}
%\usepackage{cite}
%\usepackage[pdftex]{graphicx}
\usepackage{balance}
\usepackage{hhline}
\usepackage{float}
\usepackage{xcolor}



%%%%%%%%%%%%%%%%%%%%%%%%%%%%%%%%%%%%%
%% DO NOT DELETE!!
%%%%%%%%%%%%%%%%%%%%%%%%%%%%%%%%%%%%%s
%\usepackage{tikz}
%\usetikzlibrary{trees}

\usepackage{epsfig}
\usepackage{multirow}
\usepackage{url}

%%%%%%%%%%%%%%%%%%%%%%%%%%%%%%%%%%%%%%%%%%
% Enumerate and Itemize modifications
%\let\labelindent\relax

%\usepackage{enumitem}

%\sloppy

%\setlist{topsep=0pt,noitemsep} \setitemize[1]{label=$\circ$}
%%%%%%%%%%%%%%%%%%%%%%%%%%%%%%%%%%%%%%%%%%%



% The \icmltitle you define below is probably too long as a header.
% Therefore, a short form for the running title is supplied here:
%\icmltitlerunning{Submission and Formatting Instructions for ICML 2025}
%%%%%%%%%%%%%%%%%%%%%%%%%%%%%%%%%%%%%%%%%%
% Enumerate and Itemize modifications
\let\labelindent\relax


\sloppy




\newcommand{\rtable}[1]{\ensuremath{\mathsf{#1}}}
\newcommand{\ratt}[1]{\ensuremath{\mathit{#1}}}
\newcommand{\at}[1]{\protect\ensuremath{\mathsf{#1}}\xspace}
\newcommand{\myhrule}{\rule[.5pt]{\hsize}{.5pt}}
\newcommand{\oneurl}[1]{\texttt{#1}}
\newcommand{\eat}[1]{}
\newcommand{\stab}{\rule{0pt}{8pt}\\[-1.2ex]}
\newcommand{\sttab}{\rule{0pt}{8pt}\\[-2.7ex]}
\newcommand{\sstab}{\rule{0pt}{8pt}\\[-1.8ex]}
\newcommand{\tabstrut}{\rule{0pt}{4pt}\vspace{-0.07in}}
\newcommand{\vs}{\vspace{1ex}}
\newcommand{\exa}[2]{{\tt\begin{tabbing}\hspace{#1}\=\+\kill #2\end{tabbing}}}
\newcommand{\ra}{\rightarrow}
\newcommand{\Ra}{\Rightarrow}
\newcommand{\la}{\leftarrow}
\newcommand{\bi}{\begin{itemize}}
	\newcommand{\ei}{\end{itemize}}
\newenvironment{tbi}{\begin{itemize}
		\setlength{\topsep}{0.6ex}\setlength{\itemsep}{0ex}} %\vspace{-0.5ex}}
{\end{itemize}} %\vspace{-0.5ex}}
\newenvironment{tbe}{\begin{enumerate}
\setlength{\topsep}{0ex}\setlength{\itemsep}{0ex}\vspace{0ex}}
{\end{itemize}\vspace{-1ex}}


\newcommand{\mat}[2]{{\begin{tabbing}\hspace{#1}\=\+\kill #2\end{tabbing}}}
\newcommand{\m}{\hspace{0.05in}}
\newcommand{\ls}{\hspace{0.1in}}
\newcommand{\be}{\begin{enumerate}}
\newcommand{\ee}{\end{enumerate}}
%\newcommand{\beqn}{\begin{eqnarray*}}
%\newcommand{\eeqn}{\end{eqnarray*}}
\newcommand{\beqn}{\begin{eqnarray}}
\newcommand{\eeqn}{\end{eqnarray}}

\newcommand{\card}[1]{\mid\! #1\!\mid}
\newcommand{\fth}{\hfill $\Box$}
\newcommand{\AND}{\displaystyle{\bigwedge_{i=1}^{n}}}
%\newcommand{\U}[1]{\displaystyle{\bigcup_{#1}}}
\newcommand{\Sm}[1]{\displaystyle{\sum_{#1}}}
\newcommand{\stitle}[1]{\vspace{1.6ex}\noindent{\bf #1}}
\newcommand{\sstitle}[1]{\vspace{1ex}\noindent{\bf #1}}
\newcommand{\etitle}[1]{\vspace{1ex}\noindent{\underline{\em #1}}}
\newcommand{\eetitle}[1]{\vspace{0.8ex}\noindent{{\em #1}}}
\renewcommand{\t}{\tau}
\newcommand{\Inh}[1]{\$#1}
\renewcommand{\r}[1]{{\it rule}(#1)}
\newcommand{\pa}{\parallel}
\newcommand{\LHS}{\kw{{\small LHS}}\xspace}
\newcommand{\RHS}{\kw{RHS}\xspace}
\newcommand{\ie}{\emph{i.e.,}\xspace}
\newcommand{\eg}{\emph{e.g.,}\xspace}
\newcommand{\wrt}{\emph{w.r.t.}\xspace}
\newcommand{\aka}{\emph{a.k.a.}\xspace}
\newcommand{\kwlog}{\emph{w.l.o.g.}\xspace}
\newcommand{\reffig}[1]{Fig.\ref{#1}}
\newcommand{\refalg}[1]{Algorithm \ref{#1}}
\newcommand{\refsec}[1]{Section \ref{#1}}
\newcommand{\refeq}[1]{Eq \ref{#1}}
\newcommand{\reftab}[1]{Table \ref{#1}}
\newcommand{\refprf}[1]{Prf. \ref{#1}}
%%%%%%%%%%%%%%%%%%%%%%%%%%%%%%%%%%%%%%%%%%%%%%%%%%%%%%%%%%%%%%%%
%                  Relation Algebra operators
%%%%%%%%%%%%%%%%%%%%%%%%%%%%%%%%%%%%%%%%%%%%%%%%%%%%%%%%%%%%%%%%

\newcommand{\RS}{{\small S}\xspace}
\newcommand{\RP}{{\small P}\xspace}
\newcommand{\RJ}{{\sc j}\xspace}
\newcommand{\RC}{{\small C}\xspace}
\newcommand{\RSJ}{{\small SJ}\xspace}
\newcommand{\RSC}{{\small SC}\xspace}
\newcommand{\RSP}{{\small SP}\xspace}
\newcommand{\RPJ}{{\small PJ}\xspace}
\newcommand{\RPC}{{\small PC}\xspace}
\newcommand{\RSPJ}{{\sc spj}\xspace}
\newcommand{\RSPC}{{\small SPC}\xspace}
\newcommand{\RSPJU}{{\sc spju}\xspace}
\newcommand{\RSPCU}{{\small SPCU}\xspace}
\newcommand{\RSPJUN}{{\small SPJU$^N$}\xspace}
\newcommand{\RSPCUN}{{\small SPCU$^N$}\xspace}
%%%%%%%%%%%%%%%%%%%%%%%%%%%%%%%%%%%%%%%%%%%%%%%%%%%%%%%%%%%%%%%%%%%%%%%%%%%%%%
% ALGORITHMS
%%%%%%%%%%%%%%%%%%%%%%%%%%%%%%%%%%%%%%%%%%%%%%%%%%%%%%%%%%%%%%%%%%%%%%%%%%%%%%%
\newcommand{\SELECT}{\mbox{{\bf select}}\ }
\newcommand{\FROM}{\mbox{{\bf from}\ }}
\newcommand{\WHERE}{\mbox{\bf where}\ }
\newcommand{\SUM}{\mbox{{\bf sum}}\ }
\newcommand{\GROUPBY}{\mbox{{\bf group by}}\ }
\newcommand{\HAVING}{\mbox{{\bf having}}\ }
\newcommand{\CASE}{\mbox{{\bf case}}\ }
\newcommand{\END}{\mbox{{\bf end}}\ }
\newcommand{\WHEN}{\mbox{{\bf when}}\ }
\newcommand{\EXISTS}{\mbox{{\bf exists}}\ }
\newcommand{\COUNT}{\mbox{\kw{count}}}
\newcommand{\INSERTINTO}{\mbox{{\bf insert into}}\ }
\newcommand{\UPDATE}{\mbox{{\bf update}}\ }
\newcommand{\SET}{\mbox{{\bf set}}\ }
\newcommand{\IN}{\mbox{{\bf in}}\ }
\newcommand{\If}{\mbox{\bf if}\ }
\newcommand{\Let}{\mbox{\bf let}\ }
\newcommand{\Call}{\mbox{\bf call}\ }
\newcommand{\Then}{\mbox{\bf then}\ }
\newcommand{\To}{\mbox{\bf to}\ }
\newcommand{\Else}{\mbox{\bf else}\ }
\newcommand{\ElseIf}{\mbox{\bf elseif}\ }
\newcommand{\While}{\mbox{\bf while}\ }
\newcommand{\Begin}{\mbox{\bf begin}\ }
\newcommand{\End}{\mbox{\bf end}\ }
\newcommand{\Do}{\mbox{\bf do}\ }
\newcommand{\Downto}{\mbox{\bf downto}\ }
\newcommand{\Repeat}{\mbox{\bf repeat}\ }
\newcommand{\Until}{\mbox{\bf until}\ }
\newcommand{\For}{\mbox{\bf for}\ }
\newcommand{\Each}{\mbox{\bf each}\ }

\newcommand{\ForEach}{\mbox{\bf for each}\ }
\newcommand{\Or}{\mbox{\bf or}\ }
\newcommand{\Not}{\mbox{\bf not}\ }
\newcommand{\Break}{\mbox{\bf break}\ }
\newcommand{\Continue}{\mbox{\bf continue}\ }
\newcommand{\Return}{\mbox{\bf return}\ }
\newcommand{\Case}{\mbox{\bf case}\ }
\newcommand{\Of}{\mbox{\bf of}\ }
\newcommand{\EndCase}{\mbox{\bf end-case}\ }
\newcommand{\NIL}{\mbox{\em nil}}
\newcommand{\False}{\mbox{\em false}\xspace}
\newcommand{\True}{\mbox{\em true}\xspace}
\newcommand{\algAND}{{\sc and}\xspace}
\newcommand{\OR}{{\sc or}\xspace}
\newcommand{\NOT}{{\sc not}\xspace}
\newcommand{\kw}[1]{{\ensuremath {\mathsf{#1}}}\xspace}
\newcommand{\Reps}{S}
\newcounter{ccc}
\newcommand{\bcc}{\setcounter{ccc}{1}\theccc.}
\newcommand{\bccc}{\theccc.}
\newcommand{\icc}{\addtocounter{ccc}{1}\theccc.}
\newcommand{\checking}{{\mbox{\small\sf Checking}\xspace}}
\newcommand{\preProcessing}{{\mbox{\small\sf preProcessing}\xspace}}
\newcommand{\CFDconsistency}{{\mbox{\small\sf CFD\_Checking}\xspace}}
\newcommand{\MCS} {\kw{MCS}}
\newcommand{\EIP} {\kw{EIP}}
\newcommand{\CDP} {\kw{DMP}}
\newcommand{\templateDB}{{\mbox{\small\sf templateDB}\xspace}}
\newcommand{\ChaseChecking}{{\mbox{\small\sf RandomChecking}\xspace}}
\newcommand{\chase}{{\mbox{\small\sf Chase}\xspace}}
\newcommand{\SAT}{{\mbox{\small\sf SAT}\xspace}}
\newcommand{\kSAT}{{\mbox{\small 3SAT}\xspace}}
\newcommand{\PropCFDSPC}{\kw{Prop{\small CFD\_SPC}}}
\newcommand{\PropCFDSPCU}{\kw{Prop{\small CFD\_SPCU}}}
\newcommand{\UnionEQs}{\kw{UnionEQs}}
\newcommand{\UnionCFDs}{\kw{UnionCFDs}}
\newcommand{\EQ}{\kw{EQ}}
\newcommand{\eq}{\kw{eq}}
\newcommand{\key}{\kw{key}}
\newcommand{\rep}{\kw{rep}}
\newcommand{\PEQ}{\kw{EQ2CFD}}
\newcommand{\Drop}{\kw{Drop}}
\newcommand{\Res}{\kw{Res}}
\newcommand{\CFD}{\kw{CFD}}
\newcommand{\CFDs}{\kw{CFDs}}
\newcommand{\MD}{\kw{MD}}
\newcommand{\MDs}{\kw{MDs}}
\newcommand{\CMDs}{\kw{CMDs}}
\newcommand{\MRML}{\kw{MRL}}
\newcommand{\MRMLs}{\kw{MRLs}}

\newcommand{\DC}{\kw{DC}}
\newcommand{\DCs}{\kw{DCs}}
\newcommand{\CIND}{{\sc cind}\xspace}
\newcommand{\cind}{{\small \sf CIND}}
\newcommand{\cfd}{{\small \sf CFD}}
\newcommand{\CINDp}{{\sc cind}$^+$\xspace}
\newcommand{\CINDn}{{\sc cind}$^-$\xspace}
\newcommand{\CINDs}{{\sc cind}{\small s}\xspace}
\newcommand{\FD}{{\small FD}\xspace}
\newcommand{\FDs}{\kw{FDs}}
\newcommand{\IND}{{\sc ind}\xspace}
\newcommand{\INDs}{{\sc ind}{\small s}\xspace}
\newcommand{\TGDs}{{\sc tgd}{\small s}\xspace}
\newcommand{\NP}{\kw{NP}}
\newcommand{\DAGs}{{\small DAG}s\xspace}
\newcommand{\NC}{{\sc nc}\xspace}
\newcommand{\coNP}{\kw{coNP}}
\newcommand{\PTIME}{\kw{PTIME}}
\newcommand{\PSPACE}{{\sc pspace}\xspace}
\newcommand{\EXPTIME}{{\sc exptime}\xspace}
\newcommand{\NPSPACE}{{\sc npspace}\xspace}
\newcommand{\dom}{\protect\ensuremath{\mathsf{dom}}\xspace}
\newcommand{\atset}{\protect\ensuremath{\mathsf{attr}}\xspace}
\newcommand{\attr}[1]{\protect\ensuremath{\mathsf{#1}}\xspace}
\newcommand{\attrset}{\protect\ensuremath{\mathsf{attr}}\xspace}
\newcommand{\finatset}{\protect\ensuremath{\mathsf{finattr}}\xspace}
\newcommand{\pvar}{\protect\ensuremath{\mathsf{var\%}}\xspace}
\newcommand{\lLHS}{\protect\ensuremath{\mathsf{{\small LHS}}}\xspace}
\newcommand{\RA}{{\small RA}\xspace}
\newcommand{\RBR}{\kw{RBR}}
\newcommand{\SQL}{\kw{SQL}}
\newcommand{\XSLT}{{\sc xslt}\xspace}
\newcommand{\DBMS}{{\sc dbms}\xspace}
\newcommand{\ATG}{{\sc atg}\xspace}
\newcommand{\ATGs}{{\sc atg}{\small s}\xspace}
\newcommand{\EBI}{{\sc ebi}\xspace}
\newcommand{\GO}{{\sc go}\xspace}
\newcommand{\VEC}[1]{{\sc vec}(#1)}
\newcommand{\DAG}{{\small DAG}\xspace}
\newcommand{\XQ}{{\sc xq}\xspace}
\newcommand{\XQwc}{{\sc xq}$^{\scriptscriptstyle[*]}$\xspace}
\newcommand{\XQdes}{{\sc xq}$^{\scriptscriptstyle[//]}$\xspace}
\newcommand{\XQfull}{{\sc xq}$^{\scriptscriptstyle[*,//]}$\xspace}
\newcommand{\vect}[1]{$\langle$ #1 $\rangle$}
\newcommand{\sem}[1]{[\![#1]\!]}
%\newcommand{\NN}[2]{#1\sem{#2}}
\newcommand{\e}[2]{{\mathit (#1,#2)}}
\newcommand{\ep}[2]{{\mathit (#1,#2)+}}
\newcommand{\brname}{\ensuremath{{\mathsf{N}}}}
\newcommand{\budrel}[1]{\ensuremath{{\brname_{#1}}}}
\newcommand{\budgen}[2]{\ensuremath{Q^\brname_\e{#1}{#2}}}
\newcommand{\budcut}[2]{\ensuremath{Q_\e{#1}{#2}}}
\newcommand{\R}{{\mathcal R}}
\newcommand{\A}{{\mathcal A}}
\newcommand{\AD}{{\mathcal A}_\Delta}
\newcommand{\B}{{\mathcal B}}
\newcommand{\Buff}[1]{{\mathbb{B}_{#1}}}
%\newcommand{\buf}[1]{{\mathbb{B}^b_{#1}}}
\newcommand{\Q}{{\mathcal Q}}
\newcommand{\Y}{{\mathcal Y}}
\newcommand{\I}{{\mathcal I}}
\newcommand{\V}{{\mathcal V}}
\newcommand{\E}{{\mathcal E}}
\newcommand{\Lq}{{\mathcal L}}
\renewcommand{\H}{{\mathcal H}}

\newcommand{\SC}{\kw{SC}}

\newcommand{\flag}{\kw{flag}}
\newcommand{\BF}{\kw{BF}}
\newcommand{\precision}{\kw{precision}}
\newcommand{\recall}{\kw{recall}}
\newcommand{\accuracy}{\kw{accuracy}}
\newcommand{\Fm}{\kw{F}-\kw{measure}}

\newcommand{\eop}{\hspace*{\fill}\mbox{$\Box$}}     % End of proof


\newcounter{example}%[section]
\renewcommand{\theexample}{\arabic{example}}
\newenvironment{example}{
\vspace{1.5ex}
\refstepcounter{example}
{\noindent\bf Example \theexample:}}{
\eop\vspace{1.5ex}}

\def\copyrightspace{}
\renewcommand{\ni}{\noindent}
\newcommand{\comlore}[1]{\begin{minipage}{3in}\fbox{\fbox{\parbox[t]{3in}{{\vspace{2mm}\noindent \bf COMM(LORE):~
				{ #1}\hfill  END.}}}}\end{minipage}\\}
\newcommand{\comwenfei}[1]{\begin{minipage}{3in}\fbox{\fbox{\parbox[t]{3in}{{\vspace{2mm}\noindent \bf COMM(WENFEI):~
				{ #1}\hfill  END.}}}}\end{minipage}\\}
\newcommand{\comshuai}[1]{\begin{minipage}{3in}\fbox{\fbox{\parbox[t]{3in}{{\vspace{2mm}\noindent \bf COMM(SHUAI):~
				{ #1}\hfill  END.}}}}\end{minipage}\\}
\newcommand{\nthesection}{\arabic{section}}

\newcounter{theorem}
\renewcommand{\thetheorem}{\arabic{theorem}}
\newcounter{prop}
\renewcommand{\theprop}{\arabic{theorem}}
\newcounter{lemma}
\renewcommand{\thelemma}{\arabic{theorem}}
\newcounter{cor}
\renewcommand{\thecor}{\arabic{theorem}}
\newenvironment{theorem}{\begin{em}
\refstepcounter{theorem}
{\vspace{1ex} \noindent\bf  Theorem  \thetheorem:}}{
\end{em}\eop\vspace{1ex}} %\hspace*{\fill}\vspace*{1ex}}
\newenvironment{ctheorem}[1]{\begin{em}
\refstepcounter{theorem}
{\vspace{1ex}{\noindent \bf Theorem  \thetheorem~[#1]}: }}{
\end{em}\eop\vspace{1.5ex}}
%\hspace*{\fill}\mbox{$\Box$}\vspace*{1.5ex}}
\newenvironment{prop}{\begin{em}
\refstepcounter{theorem}
{\vspace{1.5ex}\noindent \bf Proposition \theprop:}}{
\end{em}\eop\vspace{1.5ex}}%\hspace*{\fill}\vspace*{1ex}}
\newenvironment{lemma}{\begin{em}
\refstepcounter{theorem}
{\vspace{1ex}\noindent\bf Lemma \thelemma:}}{
\end{em}\eop\vspace{1ex}} %\hspace*{\fill}\vspace*{1ex}}
\newenvironment{cor}{\begin{em}
\refstepcounter{theorem}
{\vspace{1.5ex}\noindent\bf Corollary \thecor:}}{
\end{em}\eop\vspace{1.5ex}} %\hspace*{\fill}\vspace*{1ex}}


\newcounter{definition}
\renewcommand{\thedefinition}{\arabic{definition}}
\newenvironment{definition}{
	\vspace{1.5ex}
	\refstepcounter{definition}
	{\noindent\bf Definition {\bf \thedefinition}:}}{\eop%\vspace{1.5ex}
}


\newcounter{alg}[section]
\renewcommand{\thealg}{\nthesection.\arabic{alg}}
\newenvironment{alg}[1]{
\refstepcounter{alg}
{\vspace{1ex}\noindent\bf Algorithm \thealg:\, #1}}{
\vspace*{1ex}}
\newcounter{arule}
\renewcommand{\thearule}{\arabic{arule}}
\newenvironment{arule}{
\vspace{0.6ex}
\refstepcounter{arule}
{\noindent \em Rule \thearule:}}{
}
\newcounter{claim}
\renewcommand{\theclaim}{\arabic{claim}}
\newenvironment{claim}{
\vspace{0.6ex}
\refstepcounter{claim}
{\noindent\em Claim \theclaim:}}{{ Wenfei Fan}\\
}
\newenvironment{proofS}{
\vspace{1ex}
{\noindent\bf Proof:\ }}{\eop\vspace{1ex}}


%\newcommand{\SIM}{\ensuremath{\mathsf{mat}}}
\newcommand{\SIM}{\kw{SIM}}

\newcommand{\SIMe}{\ensuremath{\mathsf{mat_e}}}
\renewcommand{\S}{{\mathcal E}}

\newcommand{\GPAR}{\kw{GPAR}}
\newcommand{\GPARs}{\kw{GPARs}}
\newcommand{\GRAPE}{\kw{GRAPE}}
\newcommand{\AGRAPE}{\kw{GRAPE+}}

\newcommand{\Deep}{\kw{DEEP}}

\newcommand{\suport}{\kw{supp}}
\newcommand{\confd}{\kw{conf}}
\newcommand{\cd}[1]{|\!|{#1}|\!|}

\newcommand{\tbf}{\textbf{\textcolor{red}{X}}\xspace}
\newcommand{\qual}{\kw{qual}}
\newcommand{\MaxCard}{\kw{qualCard}}
\newcommand{\MaxSim}{\kw{qualSim}}

\newcommand{\Rees}{R_{(e,e)}}
%\newcommand{\Reps}{R_{(e,p)}}
%\newcommand{\Rpps}{R_{(p,p)}}
\newcommand{\ees}{\preceq_{(e,e)}}
%\newcommand{\eps}{\preceq_{(e,p)}}
\newcommand{\nees}{\not\preceq_{e,e}}
%\newcommand{\neps}{\not\preceq_{e,p}}
%\newcommand{\pps}{\preceq_{(p,p)}}
\newcommand{\bees}{\preceq^\mathsf{bw}_{(e,e)}}
\newcommand{\nbees}{\not\preceq^\mathsf{bw}_{(e,e)}}
%\newcommand{\beps}{\preceq^\mathsf{bw}_{(e,p)}}
%\newcommand{\nbeps}{\not\preceq^\mathsf{bw}_{(e,p)}}

\newcommand{\iees}{\precsim^{1-1}_{(e,e)}}
\newcommand{\ieps}{\precsim^{1-1}_{(e,p)}}
\newcommand{\sees}{\precsim_{(e,e)}}
%\newcommand{\seps}{\precsim^s_{(e,p)}}
\newcommand{\seps}{\precsim_{(e,p)}}
\newcommand{\nseps}{\not\precsim_{(e,p)}}
\newcommand{\nieps}{\not\precsim^{1-1}_{(e,p)}}
\newcommand{\oees}{\precsim^{\kw{opt}}_{(e,e)}}
\newcommand{\oeps}{\precsim^{\kw{opt}}_{(e,p)}}
\newcommand{\bsees}{\precsim^\mathsf{bw}_{(e,e)}}
\newcommand{\bseps}{\precsim^\mathsf{bw}_{(e,p)}}
\newcommand{\compPSim}{\kw{compPSim}}

%\newcommand{\URL} {\kw{URL}}
%\newcommand{\URLs} {\kw{URLs}}
\newcommand{\WIS} {\kw{WIS}}
\newcommand{\IS} {\kw{IS}}
\newcommand{\AFPR}{\kw{AFP}-\kw{reduction}}
\newcommand{\AFPRs}{\kw{AFP}-\kw{reductions}}

\newcommand{\OPT}{\kw{opt}}
\newcommand{\obj}{\kw{obj}}
\newcommand{\N} {{\mathcal N}}
\newcommand{\kL} {{\cal L}}

\newcommand{\maxCSPS} {\kw{compMaxCard}}
\newcommand{\maxCSPI} {\kw{compMaxCard^{1-1}}}
\newcommand{\maxSSPS} {\kw{compMaxSim}}
\newcommand{\maxSSPI} {\kw{compMaxSim^{1-1}}}
\newcommand{\subIso} {\kw{cdkMCS}}
\newcommand{\combine} {\kw{combinedMaxSim}}
\newcommand{\SF}{\kw{SF}}
%\newcommand{\VSM}{\kw{VSM}}

\newcommand{\REE}{\kw{REE}}
\newcommand{\REEs}{\kw{REEs}}

%\newcommand{\pSim}{\kw{compSimilarity}}
\newcommand{\gSim}{\kw{graphSimulation}}

\newcommand{\maxWIS} {\kw{compMaxWIS}}
\newcommand{\nei} {\kw{Neighbor}}
\newcommand{\nonNei} {\kw{NonNeighbor}}
\newcommand{\ramsey} {\kw{Ramsey}}
\newcommand{\cRamsey} {\kw{ISRemoval}}
\newcommand{\wis} {\kw{maxWIS}}
\newcommand{\maxWeight} {\kw{maxWeight}}

\newcommand{\naive} {\kw{Naive}}
\newcommand{\good} {\kw{good}}
\newcommand{\bad} {\kw{minus}}

\newcommand{\static} {\kw{static}}
\newcommand{\parent} {\kw{prev}}
\newcommand{\child} {\kw{post}}
\newcommand{\greedy} {\kw{greedyMatch}}
\newcommand{\proNeighbor} {\kw{trimMatching}}

\newcommand{\Pick}{\mbox{\bf pick}\ }

\newcommand{\sizeof} {\kw{sizeof}}
\newcommand{\genPG} {\kw{genPGraph}}
\newcommand{\genSOL} {\kw{genSolution}}

\renewcommand{\texttt}[1]{{\small\textsf{#1}}}
\newcommand{\kchanged}{\kw{\small CHANGED}}

\newcommand{\added}[1]{\textcolor{blue}{#1}}
\newcommand{\changed}[1]{\textcolor{red}{#1}}
\newcommand{\removed}[1]{\textcolor{gray}{#1}}
\newcommand{\APX}{\kw{APX}-\kw{hard}}
\newcommand{\WVC}{\kw{MIN}-\kw{WVCP}}
\newcommand{\WLM}{\kw{MIN}-\kw{WLMP}}
\newcommand{\Claim}{\kw{Claim}}
\newcommand{\MWLM}{\kw{Minimum Weighted Landmark}-\kw{Cover}}
%\newcommand{\problem}{\kw{Problem}}
\newcommand{\setcover}{\kw{SET}-\kw{COVER}}
\newcommand{\len}{\kw{len}}
\newcommand{\dis}{\kw{dis}}
\newcommand{\khopjoin}{\kw{k}-\kw{hop}-\kw{join}}
\newcommand{\Matrix} {\kw{Matrix}}
\newcommand{\BFS} {\kw{BFS}}
\newcommand{\BIBFS} {\kw{BIBFS}}
\newcommand{\twohop} {\kw{2}-\kw{hop}}
\newcommand{\dijkstra} {\kw{Dijkstra}}
\newcommand{\inport}{\kw{In}-\kw{Port}}
\newcommand{\outport}{\kw{Out}-\kw{Port}}

\newcommand{\conv}{\kw{conv}}
\newcommand{\lift}{\kw{lift}}
\newcommand{\supp}{\kw{supp}}
\newcommand{\join}{\kw{Join}}
\newcommand{\Union}{\bigcup}
\newcommand{\ak}{\allowbreak}

\newcommand{\MRC}{{\cal MRC}\xspace}
%\newcommand{\MR}{\kw{DMatch}}

\newcommand{\Assemble}{\kw{Assemble}}
\newcommand{\PEval}{\kw{PEval}}
\newcommand{\IncEval}{\kw{IncEval}}
\newcommand{\T}{{\mathcal T}}
\newcommand{\F}{{\mathcal F}}
\newcommand{\M}{{\mathcal M}}

\newcommand{\MLConstruct}{{\mathcal C}}
\newcommand{\D}{{\mathcal D}}
\newcommand{\Mes}{\kw{MSeg}}
\newcommand{\ID}[1]{\kw{ID}_{#1}}
\newcommand{\id}{\kw{id}}

\renewcommand{\P}{{\mathcal P}}
\newcommand{\dist}{\kw{dist}}

\newcommand{\SCC}{\kw{SCC}}
\newcommand{\SCCs}{\kw{SCCs}}
\newcommand{\CC}{\kw{CC}}
\newcommand{\CCs}{\kw{CCs}}
\newcommand{\ccs}{{\sl CC}s\xspace}
\newcommand{\cc}{{\sl CC}\xspace}
\newcommand{\scc}{\kw{scc}}
\newcommand{\MST}{\kw{MST}}
\newcommand{\SSSP}{\kw{SSSP}}
\newcommand{\kNN}{\kw{kNN}}
\newcommand{\knn}{\kw{knn}}
\newcommand{\rvec}{\kw{vec}}
\newcommand{\true}{\kw{true}}
\newcommand{\false}{\kw{false}}

\newcommand{\dfs}{\kw{DFS}}
\newcommand{\aff}{\kw{AFF}}

\newcommand{\spark}{\kw{Spark}}

%\newcommand{\revise}[1]{#1}
\newcommand{\revise}[1]{{\color{black}{#1}}}
%\newcommand{\revise}[1]{{\color{blue}{#1}}}
\newcommand{\revisef}[1]{{\color{black}{#1}}}
\newcommand{\reviserdy}[1]{{\color{blue}{#1}}}

%\newcommand{\change}[1]{{\underline{\textcolor{blue}{#1}}}}
%\newcommand{\change}[1]{#1}
\newcommand{\warn}[1]{{\color{black}{#1}}}
%\newcommand{\warn}[1]{{\color{red}{#1}}}
\newcommand{\warnf}[1]{{\color{black}{#1}}}
\newcommand{\blue}[1]{{\color{black}{#1}}}
\newcommand{\warnff}[1]{{\color{black}{#1}}}
\newcommand{\warnfff}[1]{{\color{black}{#1}}}
%\newcommand{\warn}[1]{#1}

%% experiments
\newcommand{\Sim}{\kw{Sim}}
\newcommand{\Sub}{\kw{SubIso}}
\newcommand{\keyword}{\kw{Keyword}}

\newcommand{\dbpedia}{\kw{DBpedia}}
\newcommand{\yago}{\kw{YAGO}}
\newcommand{\pokec}{\kw{Pokec}}
\newcommand{\live}{\kw{liveJournal}}
\newcommand{\traffic}{\kw{traffic}}
\newcommand{\uk}{\kw{UKWeb}}
\newcommand{\friend}{\kw{Friendster}}
\newcommand{\netflix}{\kw{Netflix}}

\newcommand{\giraph}{\kw{Giraph}}
\newcommand{\glab}{\kw{GraphLab}}
\newcommand{\blogel}{\kw{Blogel}}
\newcommand{\maiter}{\kw{Maiter}}
\newcommand{\guc}{\kw{GiraphUC}}
\newcommand{\hsync}{\kw{PowerSwitch}}
\newcommand{\petuum}{\kw{Petuum}}


\newcommand{\grapeni}{$\kw{GRAPE_{NI}}$\xspace}
\newcommand{\gskew}{\kw{skew}}
\newcommand{\PIE}{\kw{PIE}}

\newcommand{\CF}{\kw{CF}}
\newcommand{\PR}{\kw{PageRank}}
\newcommand{\pagerank}{\kw{pageRank}}

\newcommand{\GBSP}{$\AGRAPE_{\BSP}$\xspace}
\newcommand{\GAP}{$\AGRAPE_{\AP}$\xspace}
\newcommand{\GSSP}{$\AGRAPE_{\SSP}$\xspace}
\newcommand{\glabsync}{$\glab_{\kw{sync}}$\xspace}
\newcommand{\glabasync}{$\glab_{\kw{async}}$\xspace}
\newcommand{\BSP}{\kw{BSP}}
\newcommand{\AP}{\kw{AP}}
\newcommand{\SSP}{\kw{SSP}}
\newcommand{\AAP}{\kw{AAP}}

\newcommand{\TTL}{\kw{DS}}
\newcommand{\TIMES}{\kw{X}}


%\newlist{myitemize}{itemize}{3}
%\setlist[myitemize,1]{label=$\circ$,leftmargin=3.6ex}
%\setlist[myitemize,2]{label=$\bullet$,leftmargin=4.0ex}
%\setlist[myitemize,3]{label=$\diamond$, leftmargin=3.5ex}
%\newenvironment{mbi}{\begin{myitemize}
%        \setlength{\topsep}{0.6ex}\setlength{\itemsep}{0.16ex}\vspace{0.36ex}}
%        {\end{myitemize}\vspace{0.16ex}}

%        \newcommand{\mei}{\end{myitemize}\vspace{0.6ex}}
%\newenvironment{tmbi}{\begin{myitemize}
%       \setlength{\topsep}{0.36ex}\setlength{\itemsep}{0ex}}
%       {\end{myitemize}\vspace{1ex}}


\makeatletter
\newcommand\figcaption{\def\@captype{figure}\caption}
\newcommand\tabcaption{\def\@captype{table}\caption}
\makeatother

\newcommand{\EU}{\kw{L}}

\newcommand{\ids}{\kw{ids}}

\newcommand{\EGD}{\kw{EGD}}
\newcommand{\EGDs}{\kw{EGDs}}

\newcommand{\Eq}{\kw{Eq}}
\newcommand{\Eqv}{\kw{Eq}_\varphi}
\newcommand{\NEq}{\kw{NEq}}


\newcommand{\shpacc}{\kw{shpacc}}
\newcommand{\tktacc}{\kw{tktacc}}
\newcommand{\dev}{\kw{dev}}
\newcommand{\trk}{\kw{trk}}
\newcommand{\name}{\kw{name}}
\newcommand{\City}{\kw{city}}
\newcommand{\ai}{\kw{ai}}
\newcommand{\no}{\kw{no.}}
\newcommand{\pid}{\kw{pid}}
\newcommand{\sid}{\kw{sid}}
\newcommand{\zip}{\kw{zip}}
\newcommand{\addr}{\kw{addr}}
\newcommand{\ph}{\kw{ph}}
\newcommand{\gndr}{\kw{gndr}}
\newcommand{\ppref}{\kw{ppref}}
\newcommand{\dob}{\kw{dob}}
\newcommand{\epref}{\kw{epref}}
\newcommand{\imei}{\kw{imei}}
\newcommand{\model}{\kw{model}}
\newcommand{\ipl}{\kw{ipl}}
\newcommand{\shno}{\kw{shno.}}
\newcommand{\dai}{\kw{dai}}
\newcommand{\shai}{\kw{shai}}
\newcommand{\dno}{\kw{dno.}}
\newcommand{\srch}{\kw{srch}}
\newcommand{\visit}{\kw{visit}}
\newcommand{\eid}{\kw{eid}}
\newcommand{\knull}{\kw{null}}

\newcommand{\seller}{\kw{seller}}
\newcommand{\type}{\kw{type}}
\newcommand{\fee}{\kw{fee}}
%\newcommand{\city}{\kw{city}}
\newcommand{\weight}{\kw{weight}}
\newcommand{\product}{\kw{product}}
\newcommand{\descr}{\kw{description}}
\newcommand{\quantity}{\kw{quantity}}
\newcommand{\price}{\kw{price}}
\newcommand{\online}{\kw{online\_year}}
\newcommand{\create}{\kw{date\_created}}
\newcommand{\sale}{\kw{on\_sale\_product}}
\newcommand{\shop}{\kw{shop}}
\newcommand{\items}{\kw{item}}
\newcommand{\delivery}{\kw{delivery}}

\newcommand{\CMap}{\kw{CMap}}
\newcommand{\CExtend}{\kw{CExtend}}
\newcommand{\Fix}{\kw{Fix}}
\newcommand{\Validate}{\kw{Validate}}

\newcommand{\enhance}{\kw{Enhance}}
\newcommand{\denhance}{\kw{DEnhance}}
\newcommand{\tid}{\kw{tid}}
%\newcommand{\incchase}{{\mbox{\small\sf IncChase}\xspace}\xspace}
%\newcommand{\pchase}{{\mbox{\small\sf PChase}\xspace}\xspace}

\newcommand{\incchase}{\kw{IncChase}}
\newcommand{\pchase}{\kw{PChase}}

\newcommand{\reviseW}[1]{#1}
\newcommand{\reviseWR}[1]{{\color{red}{#1}}}
\newcommand{\reviseR}[1]{#1}
\newcommand{\reviseB}[1]{#1}
%\newcommand{\cwarn}[1]{{\color{red}{#1}}}
\newcommand{\cwarn}[1]{#1}

\newcommand{\paper}{\kw{Paper}}
\newcommand{\finance}{\kw{Finance}}
\newcommand{\synthetic}{\kw{Syn}}
\newcommand{\movie}{\kw{Movie}}


\newcommand{\err}{\kw{err}}
\newcommand{\dup}{\kw{dup}}
\newcommand{\rlv}{\kw{rlv}}

\newcommand{\imdbm}{\kw{imdbm}}
\newcommand{\tmdbm}{\kw{tmdbm}}
\newcommand{\prdb}{\kw{prdb}}
\newcommand{\myear}{\kw{year}}
\newcommand{\budget}{\kw{budget}}
\newcommand{\orilan}{\kw{language}}
\newcommand{\dirct}{\kw{director}}
\newcommand{\mtitle}{\kw{title}}
\newcommand{\releasetime}{\kw{time}}
\newcommand{\length}{\kw{len}}
\newcommand{\mcountry}{\kw{country}}
\newcommand{\birthyear}{\kw{birth}}
\newcommand{\deathyear}{\kw{death}}
\newcommand{\profession}{\kw{job}}
\newcommand{\coworker}{\kw{coworker}}
\newcommand{\cast}{\kw{cast}}
\newcommand{\director}{\kw{director}}
\newcommand{\movies}{\kw{movies}}

\newcommand{\BigDansing}{\kw{BigDansing}}
\newcommand{\ProgressiveER}{\kw{ProgressiveER}}
\newcommand{\GGCP}{\kw{GGCP}}
\newcommand{\Sd}{{\mathcal S}}

\newcommand{\BCQE}{\kw{BCQE}}

\newcommand{\HTP}{\kw{HTP}}

\newcommand{\CGDs}{\kw{CGDs}}

%\newcommand{\ie}{\emph{i.e.,}\xspace}


\newcommand{\MRM}{\kw{MRLMatch}}
\newcommand{\SE}{\kw{SMatch}}
\newcommand{\PE}{\kw{PMatch}}
\newcommand{\IE}{\kw{IMatch}}

\newcommand{\Local}{\kw{LocalMatch}}
\newcommand{\TC}{\kw{Closure}}
\newcommand{\Mg}{\kw{Merge}}

\newcommand{\MHSP}{\kw{MHFP}}

\newcommand{\CQ}{\kw{CQ}}
\newcommand{\CQs}{\kw{CQs}}

\newcommand{\freq}{\kw{freq}}


\newcommand{\HQ}{\kw{HyperCube}}
\newcommand{\PT}{\kw{HyPart}}

\newcommand{\car}{\kw{Car}}

\newcommand{\HA}{\kw{HashAssign}}

\newcommand{\MN}{\kw{MergeNode}}

\newcommand{\tabincell}[2]{
\begin{tabular}{@{}#1@{}}#2\end{tabular}
}

%\floatsetup[table]{capposition=top}



\newcommand{\HC}{\kw{HC}}
\newcommand{\MQO}{\kw{MQO}}
\newcommand{\SeqM}{\kw{Match}}

\newcommand{\HGIN}{\kw{HGIN}}


\newcommand{\IDGNN}{\kw{IDGNN}}

\newcommand{\GNNs}{\kw{GNNs}}


\newcommand{\HFrame}{\kw{HFrame}}

%\newcommand{\C}{{\mathcal C}}

\newcommand{\cdd}[1]{|\!|{#1}|\!|}

\newcommand{\dual}{\kw{DualSim}}


\newcommand{\SHP}{\kw{SHP}}

\begin{document}
	

	\title{Improving Algorithms
		for Subgraph Homomorphsim via Graph Neural Networks}
		
		\maketitle
	
	% It is OKAY to include author information, even for blind
	% submissions: the style file will automatically remove it for you
	% unless you've provided the [accepted] option to the icml2025
	% package.
	
	% List of affiliations: The first argument should be a (short)
	% identifier you will use later to specify author affiliations
	% Academic affiliations should list Department, University, City, Region, Country
	% Industry affiliations should list Company, City, Region, Country
	
	% You can specify symbols, otherwise they are numbered in order.
	% Ideally, you should not use this facility. Affiliations will be numbered
	% in order of appearance and this is the preferred way.
%	\icmlsetsymbol{equal}{*}
	

	
%	\vskip 0.3in
	
	
	% this must go after the closing bracket ] following \twocolumn[ ...
	
	% This command actually creates the footnote in the first column
	% listing the affiliations and the copyright notice.
	% The command takes one argument, which is text to display at the start of the footnote.
	% The \icmlEqualContribution command is standard text for equal contribution.
	% Remove it (just {}) if you do not need this facility.
	
	%\printAffiliationsAndNotice{}  % leave blank if no need to mention equal contribution
%	\printAffiliationsAndNotice{\icmlEqualContribution} % otherwise use the standard text.





\begin{abstract}
Homomorphism is an important
structure-preserving mapping between two graphs,
and is widely used in constraint satisfaction problems, 
social network analysis, and graph databases querying et. al.
Given a graph $G$ and a pattern $Q$, 
the subgraph homomorphism problem is to decide whether there exists
a mapping $h$ from $Q$ to $G$ such that 
if there exists an edge from $u$ to $v$ in $Q$,
$G$ also has an edge from $h(u)$ to $h(v)$.
Different from the classic isomorphic mapping 
that is injective, 
homomorphism %maps 
\revise{may map} two \revise{distinct} vertices in $Q$
to the same vertex in $G$, 
and is more challenging than isomorphic mapping.
There exist multiple Graph Neural Netowrks (GNNs)
for (subgraph) ismomorphism, 
but few models are developed for subgraph-homomorphism.
We develop the first GNN-based framework 
for subgraph homomorphism. We show that 
the framework is more expressive than the vanilla GNN,
\ie it can distinguish more graphs that are 
not homomorphic.
%Moreover, we improve the performance of algorithms
%by plugging in the extended models.
We also a generalization bound for the proposed framework.
Using real-life and synthetic graphs, we 
show that our framework is faster for subgraph matching,
and its accuracy is also high.
\end{abstract}
%\textit{}

%\newcommand{\HGIN}{\kw{HGIN}}


%\eat{
%%%%%%%%%%%%%%%%%%%%%%%%%%%%%%%%%%%%%%%%%%%
% Enumerate and Itemize modifications
\let\labelindent\relax


\sloppy




\newcommand{\rtable}[1]{\ensuremath{\mathsf{#1}}}
\newcommand{\ratt}[1]{\ensuremath{\mathit{#1}}}
\newcommand{\at}[1]{\protect\ensuremath{\mathsf{#1}}\xspace}
\newcommand{\myhrule}{\rule[.5pt]{\hsize}{.5pt}}
\newcommand{\oneurl}[1]{\texttt{#1}}
\newcommand{\eat}[1]{}
\newcommand{\stab}{\rule{0pt}{8pt}\\[-1.2ex]}
\newcommand{\sttab}{\rule{0pt}{8pt}\\[-2.7ex]}
\newcommand{\sstab}{\rule{0pt}{8pt}\\[-1.8ex]}
\newcommand{\tabstrut}{\rule{0pt}{4pt}\vspace{-0.07in}}
\newcommand{\vs}{\vspace{1ex}}
\newcommand{\exa}[2]{{\tt\begin{tabbing}\hspace{#1}\=\+\kill #2\end{tabbing}}}
\newcommand{\ra}{\rightarrow}
\newcommand{\Ra}{\Rightarrow}
\newcommand{\la}{\leftarrow}
\newcommand{\bi}{\begin{itemize}}
	\newcommand{\ei}{\end{itemize}}
\newenvironment{tbi}{\begin{itemize}
		\setlength{\topsep}{0.6ex}\setlength{\itemsep}{0ex}} %\vspace{-0.5ex}}
{\end{itemize}} %\vspace{-0.5ex}}
\newenvironment{tbe}{\begin{enumerate}
\setlength{\topsep}{0ex}\setlength{\itemsep}{0ex}\vspace{0ex}}
{\end{itemize}\vspace{-1ex}}


\newcommand{\mat}[2]{{\begin{tabbing}\hspace{#1}\=\+\kill #2\end{tabbing}}}
\newcommand{\m}{\hspace{0.05in}}
\newcommand{\ls}{\hspace{0.1in}}
\newcommand{\be}{\begin{enumerate}}
\newcommand{\ee}{\end{enumerate}}
%\newcommand{\beqn}{\begin{eqnarray*}}
%\newcommand{\eeqn}{\end{eqnarray*}}
\newcommand{\beqn}{\begin{eqnarray}}
\newcommand{\eeqn}{\end{eqnarray}}

\newcommand{\card}[1]{\mid\! #1\!\mid}
\newcommand{\fth}{\hfill $\Box$}
\newcommand{\AND}{\displaystyle{\bigwedge_{i=1}^{n}}}
%\newcommand{\U}[1]{\displaystyle{\bigcup_{#1}}}
\newcommand{\Sm}[1]{\displaystyle{\sum_{#1}}}
\newcommand{\stitle}[1]{\vspace{1.6ex}\noindent{\bf #1}}
\newcommand{\sstitle}[1]{\vspace{1ex}\noindent{\bf #1}}
\newcommand{\etitle}[1]{\vspace{1ex}\noindent{\underline{\em #1}}}
\newcommand{\eetitle}[1]{\vspace{0.8ex}\noindent{{\em #1}}}
\renewcommand{\t}{\tau}
\newcommand{\Inh}[1]{\$#1}
\renewcommand{\r}[1]{{\it rule}(#1)}
\newcommand{\pa}{\parallel}
\newcommand{\LHS}{\kw{{\small LHS}}\xspace}
\newcommand{\RHS}{\kw{RHS}\xspace}
\newcommand{\ie}{\emph{i.e.,}\xspace}
\newcommand{\eg}{\emph{e.g.,}\xspace}
\newcommand{\wrt}{\emph{w.r.t.}\xspace}
\newcommand{\aka}{\emph{a.k.a.}\xspace}
\newcommand{\kwlog}{\emph{w.l.o.g.}\xspace}
\newcommand{\reffig}[1]{Fig.\ref{#1}}
\newcommand{\refalg}[1]{Algorithm \ref{#1}}
\newcommand{\refsec}[1]{Section \ref{#1}}
\newcommand{\refeq}[1]{Eq \ref{#1}}
\newcommand{\reftab}[1]{Table \ref{#1}}
\newcommand{\refprf}[1]{Prf. \ref{#1}}
%%%%%%%%%%%%%%%%%%%%%%%%%%%%%%%%%%%%%%%%%%%%%%%%%%%%%%%%%%%%%%%%
%                  Relation Algebra operators
%%%%%%%%%%%%%%%%%%%%%%%%%%%%%%%%%%%%%%%%%%%%%%%%%%%%%%%%%%%%%%%%

\newcommand{\RS}{{\small S}\xspace}
\newcommand{\RP}{{\small P}\xspace}
\newcommand{\RJ}{{\sc j}\xspace}
\newcommand{\RC}{{\small C}\xspace}
\newcommand{\RSJ}{{\small SJ}\xspace}
\newcommand{\RSC}{{\small SC}\xspace}
\newcommand{\RSP}{{\small SP}\xspace}
\newcommand{\RPJ}{{\small PJ}\xspace}
\newcommand{\RPC}{{\small PC}\xspace}
\newcommand{\RSPJ}{{\sc spj}\xspace}
\newcommand{\RSPC}{{\small SPC}\xspace}
\newcommand{\RSPJU}{{\sc spju}\xspace}
\newcommand{\RSPCU}{{\small SPCU}\xspace}
\newcommand{\RSPJUN}{{\small SPJU$^N$}\xspace}
\newcommand{\RSPCUN}{{\small SPCU$^N$}\xspace}
%%%%%%%%%%%%%%%%%%%%%%%%%%%%%%%%%%%%%%%%%%%%%%%%%%%%%%%%%%%%%%%%%%%%%%%%%%%%%%
% ALGORITHMS
%%%%%%%%%%%%%%%%%%%%%%%%%%%%%%%%%%%%%%%%%%%%%%%%%%%%%%%%%%%%%%%%%%%%%%%%%%%%%%%
\newcommand{\SELECT}{\mbox{{\bf select}}\ }
\newcommand{\FROM}{\mbox{{\bf from}\ }}
\newcommand{\WHERE}{\mbox{\bf where}\ }
\newcommand{\SUM}{\mbox{{\bf sum}}\ }
\newcommand{\GROUPBY}{\mbox{{\bf group by}}\ }
\newcommand{\HAVING}{\mbox{{\bf having}}\ }
\newcommand{\CASE}{\mbox{{\bf case}}\ }
\newcommand{\END}{\mbox{{\bf end}}\ }
\newcommand{\WHEN}{\mbox{{\bf when}}\ }
\newcommand{\EXISTS}{\mbox{{\bf exists}}\ }
\newcommand{\COUNT}{\mbox{\kw{count}}}
\newcommand{\INSERTINTO}{\mbox{{\bf insert into}}\ }
\newcommand{\UPDATE}{\mbox{{\bf update}}\ }
\newcommand{\SET}{\mbox{{\bf set}}\ }
\newcommand{\IN}{\mbox{{\bf in}}\ }
\newcommand{\If}{\mbox{\bf if}\ }
\newcommand{\Let}{\mbox{\bf let}\ }
\newcommand{\Call}{\mbox{\bf call}\ }
\newcommand{\Then}{\mbox{\bf then}\ }
\newcommand{\To}{\mbox{\bf to}\ }
\newcommand{\Else}{\mbox{\bf else}\ }
\newcommand{\ElseIf}{\mbox{\bf elseif}\ }
\newcommand{\While}{\mbox{\bf while}\ }
\newcommand{\Begin}{\mbox{\bf begin}\ }
\newcommand{\End}{\mbox{\bf end}\ }
\newcommand{\Do}{\mbox{\bf do}\ }
\newcommand{\Downto}{\mbox{\bf downto}\ }
\newcommand{\Repeat}{\mbox{\bf repeat}\ }
\newcommand{\Until}{\mbox{\bf until}\ }
\newcommand{\For}{\mbox{\bf for}\ }
\newcommand{\Each}{\mbox{\bf each}\ }

\newcommand{\ForEach}{\mbox{\bf for each}\ }
\newcommand{\Or}{\mbox{\bf or}\ }
\newcommand{\Not}{\mbox{\bf not}\ }
\newcommand{\Break}{\mbox{\bf break}\ }
\newcommand{\Continue}{\mbox{\bf continue}\ }
\newcommand{\Return}{\mbox{\bf return}\ }
\newcommand{\Case}{\mbox{\bf case}\ }
\newcommand{\Of}{\mbox{\bf of}\ }
\newcommand{\EndCase}{\mbox{\bf end-case}\ }
\newcommand{\NIL}{\mbox{\em nil}}
\newcommand{\False}{\mbox{\em false}\xspace}
\newcommand{\True}{\mbox{\em true}\xspace}
\newcommand{\algAND}{{\sc and}\xspace}
\newcommand{\OR}{{\sc or}\xspace}
\newcommand{\NOT}{{\sc not}\xspace}
\newcommand{\kw}[1]{{\ensuremath {\mathsf{#1}}}\xspace}
\newcommand{\Reps}{S}
\newcounter{ccc}
\newcommand{\bcc}{\setcounter{ccc}{1}\theccc.}
\newcommand{\bccc}{\theccc.}
\newcommand{\icc}{\addtocounter{ccc}{1}\theccc.}
\newcommand{\checking}{{\mbox{\small\sf Checking}\xspace}}
\newcommand{\preProcessing}{{\mbox{\small\sf preProcessing}\xspace}}
\newcommand{\CFDconsistency}{{\mbox{\small\sf CFD\_Checking}\xspace}}
\newcommand{\MCS} {\kw{MCS}}
\newcommand{\EIP} {\kw{EIP}}
\newcommand{\CDP} {\kw{DMP}}
\newcommand{\templateDB}{{\mbox{\small\sf templateDB}\xspace}}
\newcommand{\ChaseChecking}{{\mbox{\small\sf RandomChecking}\xspace}}
\newcommand{\chase}{{\mbox{\small\sf Chase}\xspace}}
\newcommand{\SAT}{{\mbox{\small\sf SAT}\xspace}}
\newcommand{\kSAT}{{\mbox{\small 3SAT}\xspace}}
\newcommand{\PropCFDSPC}{\kw{Prop{\small CFD\_SPC}}}
\newcommand{\PropCFDSPCU}{\kw{Prop{\small CFD\_SPCU}}}
\newcommand{\UnionEQs}{\kw{UnionEQs}}
\newcommand{\UnionCFDs}{\kw{UnionCFDs}}
\newcommand{\EQ}{\kw{EQ}}
\newcommand{\eq}{\kw{eq}}
\newcommand{\key}{\kw{key}}
\newcommand{\rep}{\kw{rep}}
\newcommand{\PEQ}{\kw{EQ2CFD}}
\newcommand{\Drop}{\kw{Drop}}
\newcommand{\Res}{\kw{Res}}
\newcommand{\CFD}{\kw{CFD}}
\newcommand{\CFDs}{\kw{CFDs}}
\newcommand{\MD}{\kw{MD}}
\newcommand{\MDs}{\kw{MDs}}
\newcommand{\CMDs}{\kw{CMDs}}
\newcommand{\MRML}{\kw{MRL}}
\newcommand{\MRMLs}{\kw{MRLs}}

\newcommand{\DC}{\kw{DC}}
\newcommand{\DCs}{\kw{DCs}}
\newcommand{\CIND}{{\sc cind}\xspace}
\newcommand{\cind}{{\small \sf CIND}}
\newcommand{\cfd}{{\small \sf CFD}}
\newcommand{\CINDp}{{\sc cind}$^+$\xspace}
\newcommand{\CINDn}{{\sc cind}$^-$\xspace}
\newcommand{\CINDs}{{\sc cind}{\small s}\xspace}
\newcommand{\FD}{{\small FD}\xspace}
\newcommand{\FDs}{\kw{FDs}}
\newcommand{\IND}{{\sc ind}\xspace}
\newcommand{\INDs}{{\sc ind}{\small s}\xspace}
\newcommand{\TGDs}{{\sc tgd}{\small s}\xspace}
\newcommand{\NP}{\kw{NP}}
\newcommand{\DAGs}{{\small DAG}s\xspace}
\newcommand{\NC}{{\sc nc}\xspace}
\newcommand{\coNP}{\kw{coNP}}
\newcommand{\PTIME}{\kw{PTIME}}
\newcommand{\PSPACE}{{\sc pspace}\xspace}
\newcommand{\EXPTIME}{{\sc exptime}\xspace}
\newcommand{\NPSPACE}{{\sc npspace}\xspace}
\newcommand{\dom}{\protect\ensuremath{\mathsf{dom}}\xspace}
\newcommand{\atset}{\protect\ensuremath{\mathsf{attr}}\xspace}
\newcommand{\attr}[1]{\protect\ensuremath{\mathsf{#1}}\xspace}
\newcommand{\attrset}{\protect\ensuremath{\mathsf{attr}}\xspace}
\newcommand{\finatset}{\protect\ensuremath{\mathsf{finattr}}\xspace}
\newcommand{\pvar}{\protect\ensuremath{\mathsf{var\%}}\xspace}
\newcommand{\lLHS}{\protect\ensuremath{\mathsf{{\small LHS}}}\xspace}
\newcommand{\RA}{{\small RA}\xspace}
\newcommand{\RBR}{\kw{RBR}}
\newcommand{\SQL}{\kw{SQL}}
\newcommand{\XSLT}{{\sc xslt}\xspace}
\newcommand{\DBMS}{{\sc dbms}\xspace}
\newcommand{\ATG}{{\sc atg}\xspace}
\newcommand{\ATGs}{{\sc atg}{\small s}\xspace}
\newcommand{\EBI}{{\sc ebi}\xspace}
\newcommand{\GO}{{\sc go}\xspace}
\newcommand{\VEC}[1]{{\sc vec}(#1)}
\newcommand{\DAG}{{\small DAG}\xspace}
\newcommand{\XQ}{{\sc xq}\xspace}
\newcommand{\XQwc}{{\sc xq}$^{\scriptscriptstyle[*]}$\xspace}
\newcommand{\XQdes}{{\sc xq}$^{\scriptscriptstyle[//]}$\xspace}
\newcommand{\XQfull}{{\sc xq}$^{\scriptscriptstyle[*,//]}$\xspace}
\newcommand{\vect}[1]{$\langle$ #1 $\rangle$}
\newcommand{\sem}[1]{[\![#1]\!]}
%\newcommand{\NN}[2]{#1\sem{#2}}
\newcommand{\e}[2]{{\mathit (#1,#2)}}
\newcommand{\ep}[2]{{\mathit (#1,#2)+}}
\newcommand{\brname}{\ensuremath{{\mathsf{N}}}}
\newcommand{\budrel}[1]{\ensuremath{{\brname_{#1}}}}
\newcommand{\budgen}[2]{\ensuremath{Q^\brname_\e{#1}{#2}}}
\newcommand{\budcut}[2]{\ensuremath{Q_\e{#1}{#2}}}
\newcommand{\R}{{\mathcal R}}
\newcommand{\A}{{\mathcal A}}
\newcommand{\AD}{{\mathcal A}_\Delta}
\newcommand{\B}{{\mathcal B}}
\newcommand{\Buff}[1]{{\mathbb{B}_{#1}}}
%\newcommand{\buf}[1]{{\mathbb{B}^b_{#1}}}
\newcommand{\Q}{{\mathcal Q}}
\newcommand{\Y}{{\mathcal Y}}
\newcommand{\I}{{\mathcal I}}
\newcommand{\V}{{\mathcal V}}
\newcommand{\E}{{\mathcal E}}
\newcommand{\Lq}{{\mathcal L}}
\renewcommand{\H}{{\mathcal H}}

\newcommand{\SC}{\kw{SC}}

\newcommand{\flag}{\kw{flag}}
\newcommand{\BF}{\kw{BF}}
\newcommand{\precision}{\kw{precision}}
\newcommand{\recall}{\kw{recall}}
\newcommand{\accuracy}{\kw{accuracy}}
\newcommand{\Fm}{\kw{F}-\kw{measure}}

\newcommand{\eop}{\hspace*{\fill}\mbox{$\Box$}}     % End of proof


\newcounter{example}%[section]
\renewcommand{\theexample}{\arabic{example}}
\newenvironment{example}{
\vspace{1.5ex}
\refstepcounter{example}
{\noindent\bf Example \theexample:}}{
\eop\vspace{1.5ex}}

\def\copyrightspace{}
\renewcommand{\ni}{\noindent}
\newcommand{\comlore}[1]{\begin{minipage}{3in}\fbox{\fbox{\parbox[t]{3in}{{\vspace{2mm}\noindent \bf COMM(LORE):~
				{ #1}\hfill  END.}}}}\end{minipage}\\}
\newcommand{\comwenfei}[1]{\begin{minipage}{3in}\fbox{\fbox{\parbox[t]{3in}{{\vspace{2mm}\noindent \bf COMM(WENFEI):~
				{ #1}\hfill  END.}}}}\end{minipage}\\}
\newcommand{\comshuai}[1]{\begin{minipage}{3in}\fbox{\fbox{\parbox[t]{3in}{{\vspace{2mm}\noindent \bf COMM(SHUAI):~
				{ #1}\hfill  END.}}}}\end{minipage}\\}
\newcommand{\nthesection}{\arabic{section}}

\newcounter{theorem}
\renewcommand{\thetheorem}{\arabic{theorem}}
\newcounter{prop}
\renewcommand{\theprop}{\arabic{theorem}}
\newcounter{lemma}
\renewcommand{\thelemma}{\arabic{theorem}}
\newcounter{cor}
\renewcommand{\thecor}{\arabic{theorem}}
\newenvironment{theorem}{\begin{em}
\refstepcounter{theorem}
{\vspace{1ex} \noindent\bf  Theorem  \thetheorem:}}{
\end{em}\eop\vspace{1ex}} %\hspace*{\fill}\vspace*{1ex}}
\newenvironment{ctheorem}[1]{\begin{em}
\refstepcounter{theorem}
{\vspace{1ex}{\noindent \bf Theorem  \thetheorem~[#1]}: }}{
\end{em}\eop\vspace{1.5ex}}
%\hspace*{\fill}\mbox{$\Box$}\vspace*{1.5ex}}
\newenvironment{prop}{\begin{em}
\refstepcounter{theorem}
{\vspace{1.5ex}\noindent \bf Proposition \theprop:}}{
\end{em}\eop\vspace{1.5ex}}%\hspace*{\fill}\vspace*{1ex}}
\newenvironment{lemma}{\begin{em}
\refstepcounter{theorem}
{\vspace{1ex}\noindent\bf Lemma \thelemma:}}{
\end{em}\eop\vspace{1ex}} %\hspace*{\fill}\vspace*{1ex}}
\newenvironment{cor}{\begin{em}
\refstepcounter{theorem}
{\vspace{1.5ex}\noindent\bf Corollary \thecor:}}{
\end{em}\eop\vspace{1.5ex}} %\hspace*{\fill}\vspace*{1ex}}


\newcounter{definition}
\renewcommand{\thedefinition}{\arabic{definition}}
\newenvironment{definition}{
	\vspace{1.5ex}
	\refstepcounter{definition}
	{\noindent\bf Definition {\bf \thedefinition}:}}{\eop%\vspace{1.5ex}
}


\newcounter{alg}[section]
\renewcommand{\thealg}{\nthesection.\arabic{alg}}
\newenvironment{alg}[1]{
\refstepcounter{alg}
{\vspace{1ex}\noindent\bf Algorithm \thealg:\, #1}}{
\vspace*{1ex}}
\newcounter{arule}
\renewcommand{\thearule}{\arabic{arule}}
\newenvironment{arule}{
\vspace{0.6ex}
\refstepcounter{arule}
{\noindent \em Rule \thearule:}}{
}
\newcounter{claim}
\renewcommand{\theclaim}{\arabic{claim}}
\newenvironment{claim}{
\vspace{0.6ex}
\refstepcounter{claim}
{\noindent\em Claim \theclaim:}}{{ Wenfei Fan}\\
}
\newenvironment{proofS}{
\vspace{1ex}
{\noindent\bf Proof:\ }}{\eop\vspace{1ex}}


%\newcommand{\SIM}{\ensuremath{\mathsf{mat}}}
\newcommand{\SIM}{\kw{SIM}}

\newcommand{\SIMe}{\ensuremath{\mathsf{mat_e}}}
\renewcommand{\S}{{\mathcal E}}

\newcommand{\GPAR}{\kw{GPAR}}
\newcommand{\GPARs}{\kw{GPARs}}
\newcommand{\GRAPE}{\kw{GRAPE}}
\newcommand{\AGRAPE}{\kw{GRAPE+}}

\newcommand{\Deep}{\kw{DEEP}}

\newcommand{\suport}{\kw{supp}}
\newcommand{\confd}{\kw{conf}}
\newcommand{\cd}[1]{|\!|{#1}|\!|}

\newcommand{\tbf}{\textbf{\textcolor{red}{X}}\xspace}
\newcommand{\qual}{\kw{qual}}
\newcommand{\MaxCard}{\kw{qualCard}}
\newcommand{\MaxSim}{\kw{qualSim}}

\newcommand{\Rees}{R_{(e,e)}}
%\newcommand{\Reps}{R_{(e,p)}}
%\newcommand{\Rpps}{R_{(p,p)}}
\newcommand{\ees}{\preceq_{(e,e)}}
%\newcommand{\eps}{\preceq_{(e,p)}}
\newcommand{\nees}{\not\preceq_{e,e}}
%\newcommand{\neps}{\not\preceq_{e,p}}
%\newcommand{\pps}{\preceq_{(p,p)}}
\newcommand{\bees}{\preceq^\mathsf{bw}_{(e,e)}}
\newcommand{\nbees}{\not\preceq^\mathsf{bw}_{(e,e)}}
%\newcommand{\beps}{\preceq^\mathsf{bw}_{(e,p)}}
%\newcommand{\nbeps}{\not\preceq^\mathsf{bw}_{(e,p)}}

\newcommand{\iees}{\precsim^{1-1}_{(e,e)}}
\newcommand{\ieps}{\precsim^{1-1}_{(e,p)}}
\newcommand{\sees}{\precsim_{(e,e)}}
%\newcommand{\seps}{\precsim^s_{(e,p)}}
\newcommand{\seps}{\precsim_{(e,p)}}
\newcommand{\nseps}{\not\precsim_{(e,p)}}
\newcommand{\nieps}{\not\precsim^{1-1}_{(e,p)}}
\newcommand{\oees}{\precsim^{\kw{opt}}_{(e,e)}}
\newcommand{\oeps}{\precsim^{\kw{opt}}_{(e,p)}}
\newcommand{\bsees}{\precsim^\mathsf{bw}_{(e,e)}}
\newcommand{\bseps}{\precsim^\mathsf{bw}_{(e,p)}}
\newcommand{\compPSim}{\kw{compPSim}}

%\newcommand{\URL} {\kw{URL}}
%\newcommand{\URLs} {\kw{URLs}}
\newcommand{\WIS} {\kw{WIS}}
\newcommand{\IS} {\kw{IS}}
\newcommand{\AFPR}{\kw{AFP}-\kw{reduction}}
\newcommand{\AFPRs}{\kw{AFP}-\kw{reductions}}

\newcommand{\OPT}{\kw{opt}}
\newcommand{\obj}{\kw{obj}}
\newcommand{\N} {{\mathcal N}}
\newcommand{\kL} {{\cal L}}

\newcommand{\maxCSPS} {\kw{compMaxCard}}
\newcommand{\maxCSPI} {\kw{compMaxCard^{1-1}}}
\newcommand{\maxSSPS} {\kw{compMaxSim}}
\newcommand{\maxSSPI} {\kw{compMaxSim^{1-1}}}
\newcommand{\subIso} {\kw{cdkMCS}}
\newcommand{\combine} {\kw{combinedMaxSim}}
\newcommand{\SF}{\kw{SF}}
%\newcommand{\VSM}{\kw{VSM}}

\newcommand{\REE}{\kw{REE}}
\newcommand{\REEs}{\kw{REEs}}

%\newcommand{\pSim}{\kw{compSimilarity}}
\newcommand{\gSim}{\kw{graphSimulation}}

\newcommand{\maxWIS} {\kw{compMaxWIS}}
\newcommand{\nei} {\kw{Neighbor}}
\newcommand{\nonNei} {\kw{NonNeighbor}}
\newcommand{\ramsey} {\kw{Ramsey}}
\newcommand{\cRamsey} {\kw{ISRemoval}}
\newcommand{\wis} {\kw{maxWIS}}
\newcommand{\maxWeight} {\kw{maxWeight}}

\newcommand{\naive} {\kw{Naive}}
\newcommand{\good} {\kw{good}}
\newcommand{\bad} {\kw{minus}}

\newcommand{\static} {\kw{static}}
\newcommand{\parent} {\kw{prev}}
\newcommand{\child} {\kw{post}}
\newcommand{\greedy} {\kw{greedyMatch}}
\newcommand{\proNeighbor} {\kw{trimMatching}}

\newcommand{\Pick}{\mbox{\bf pick}\ }

\newcommand{\sizeof} {\kw{sizeof}}
\newcommand{\genPG} {\kw{genPGraph}}
\newcommand{\genSOL} {\kw{genSolution}}

\renewcommand{\texttt}[1]{{\small\textsf{#1}}}
\newcommand{\kchanged}{\kw{\small CHANGED}}

\newcommand{\added}[1]{\textcolor{blue}{#1}}
\newcommand{\changed}[1]{\textcolor{red}{#1}}
\newcommand{\removed}[1]{\textcolor{gray}{#1}}
\newcommand{\APX}{\kw{APX}-\kw{hard}}
\newcommand{\WVC}{\kw{MIN}-\kw{WVCP}}
\newcommand{\WLM}{\kw{MIN}-\kw{WLMP}}
\newcommand{\Claim}{\kw{Claim}}
\newcommand{\MWLM}{\kw{Minimum Weighted Landmark}-\kw{Cover}}
%\newcommand{\problem}{\kw{Problem}}
\newcommand{\setcover}{\kw{SET}-\kw{COVER}}
\newcommand{\len}{\kw{len}}
\newcommand{\dis}{\kw{dis}}
\newcommand{\khopjoin}{\kw{k}-\kw{hop}-\kw{join}}
\newcommand{\Matrix} {\kw{Matrix}}
\newcommand{\BFS} {\kw{BFS}}
\newcommand{\BIBFS} {\kw{BIBFS}}
\newcommand{\twohop} {\kw{2}-\kw{hop}}
\newcommand{\dijkstra} {\kw{Dijkstra}}
\newcommand{\inport}{\kw{In}-\kw{Port}}
\newcommand{\outport}{\kw{Out}-\kw{Port}}

\newcommand{\conv}{\kw{conv}}
\newcommand{\lift}{\kw{lift}}
\newcommand{\supp}{\kw{supp}}
\newcommand{\join}{\kw{Join}}
\newcommand{\Union}{\bigcup}
\newcommand{\ak}{\allowbreak}

\newcommand{\MRC}{{\cal MRC}\xspace}
%\newcommand{\MR}{\kw{DMatch}}

\newcommand{\Assemble}{\kw{Assemble}}
\newcommand{\PEval}{\kw{PEval}}
\newcommand{\IncEval}{\kw{IncEval}}
\newcommand{\T}{{\mathcal T}}
\newcommand{\F}{{\mathcal F}}
\newcommand{\M}{{\mathcal M}}

\newcommand{\MLConstruct}{{\mathcal C}}
\newcommand{\D}{{\mathcal D}}
\newcommand{\Mes}{\kw{MSeg}}
\newcommand{\ID}[1]{\kw{ID}_{#1}}
\newcommand{\id}{\kw{id}}

\renewcommand{\P}{{\mathcal P}}
\newcommand{\dist}{\kw{dist}}

\newcommand{\SCC}{\kw{SCC}}
\newcommand{\SCCs}{\kw{SCCs}}
\newcommand{\CC}{\kw{CC}}
\newcommand{\CCs}{\kw{CCs}}
\newcommand{\ccs}{{\sl CC}s\xspace}
\newcommand{\cc}{{\sl CC}\xspace}
\newcommand{\scc}{\kw{scc}}
\newcommand{\MST}{\kw{MST}}
\newcommand{\SSSP}{\kw{SSSP}}
\newcommand{\kNN}{\kw{kNN}}
\newcommand{\knn}{\kw{knn}}
\newcommand{\rvec}{\kw{vec}}
\newcommand{\true}{\kw{true}}
\newcommand{\false}{\kw{false}}

\newcommand{\dfs}{\kw{DFS}}
\newcommand{\aff}{\kw{AFF}}

\newcommand{\spark}{\kw{Spark}}

%\newcommand{\revise}[1]{#1}
\newcommand{\revise}[1]{{\color{black}{#1}}}
%\newcommand{\revise}[1]{{\color{blue}{#1}}}
\newcommand{\revisef}[1]{{\color{black}{#1}}}
\newcommand{\reviserdy}[1]{{\color{blue}{#1}}}

%\newcommand{\change}[1]{{\underline{\textcolor{blue}{#1}}}}
%\newcommand{\change}[1]{#1}
\newcommand{\warn}[1]{{\color{black}{#1}}}
%\newcommand{\warn}[1]{{\color{red}{#1}}}
\newcommand{\warnf}[1]{{\color{black}{#1}}}
\newcommand{\blue}[1]{{\color{black}{#1}}}
\newcommand{\warnff}[1]{{\color{black}{#1}}}
\newcommand{\warnfff}[1]{{\color{black}{#1}}}
%\newcommand{\warn}[1]{#1}

%% experiments
\newcommand{\Sim}{\kw{Sim}}
\newcommand{\Sub}{\kw{SubIso}}
\newcommand{\keyword}{\kw{Keyword}}

\newcommand{\dbpedia}{\kw{DBpedia}}
\newcommand{\yago}{\kw{YAGO}}
\newcommand{\pokec}{\kw{Pokec}}
\newcommand{\live}{\kw{liveJournal}}
\newcommand{\traffic}{\kw{traffic}}
\newcommand{\uk}{\kw{UKWeb}}
\newcommand{\friend}{\kw{Friendster}}
\newcommand{\netflix}{\kw{Netflix}}

\newcommand{\giraph}{\kw{Giraph}}
\newcommand{\glab}{\kw{GraphLab}}
\newcommand{\blogel}{\kw{Blogel}}
\newcommand{\maiter}{\kw{Maiter}}
\newcommand{\guc}{\kw{GiraphUC}}
\newcommand{\hsync}{\kw{PowerSwitch}}
\newcommand{\petuum}{\kw{Petuum}}


\newcommand{\grapeni}{$\kw{GRAPE_{NI}}$\xspace}
\newcommand{\gskew}{\kw{skew}}
\newcommand{\PIE}{\kw{PIE}}

\newcommand{\CF}{\kw{CF}}
\newcommand{\PR}{\kw{PageRank}}
\newcommand{\pagerank}{\kw{pageRank}}

\newcommand{\GBSP}{$\AGRAPE_{\BSP}$\xspace}
\newcommand{\GAP}{$\AGRAPE_{\AP}$\xspace}
\newcommand{\GSSP}{$\AGRAPE_{\SSP}$\xspace}
\newcommand{\glabsync}{$\glab_{\kw{sync}}$\xspace}
\newcommand{\glabasync}{$\glab_{\kw{async}}$\xspace}
\newcommand{\BSP}{\kw{BSP}}
\newcommand{\AP}{\kw{AP}}
\newcommand{\SSP}{\kw{SSP}}
\newcommand{\AAP}{\kw{AAP}}

\newcommand{\TTL}{\kw{DS}}
\newcommand{\TIMES}{\kw{X}}


%\newlist{myitemize}{itemize}{3}
%\setlist[myitemize,1]{label=$\circ$,leftmargin=3.6ex}
%\setlist[myitemize,2]{label=$\bullet$,leftmargin=4.0ex}
%\setlist[myitemize,3]{label=$\diamond$, leftmargin=3.5ex}
%\newenvironment{mbi}{\begin{myitemize}
%        \setlength{\topsep}{0.6ex}\setlength{\itemsep}{0.16ex}\vspace{0.36ex}}
%        {\end{myitemize}\vspace{0.16ex}}

%        \newcommand{\mei}{\end{myitemize}\vspace{0.6ex}}
%\newenvironment{tmbi}{\begin{myitemize}
%       \setlength{\topsep}{0.36ex}\setlength{\itemsep}{0ex}}
%       {\end{myitemize}\vspace{1ex}}


\makeatletter
\newcommand\figcaption{\def\@captype{figure}\caption}
\newcommand\tabcaption{\def\@captype{table}\caption}
\makeatother

\newcommand{\EU}{\kw{L}}

\newcommand{\ids}{\kw{ids}}

\newcommand{\EGD}{\kw{EGD}}
\newcommand{\EGDs}{\kw{EGDs}}

\newcommand{\Eq}{\kw{Eq}}
\newcommand{\Eqv}{\kw{Eq}_\varphi}
\newcommand{\NEq}{\kw{NEq}}


\newcommand{\shpacc}{\kw{shpacc}}
\newcommand{\tktacc}{\kw{tktacc}}
\newcommand{\dev}{\kw{dev}}
\newcommand{\trk}{\kw{trk}}
\newcommand{\name}{\kw{name}}
\newcommand{\City}{\kw{city}}
\newcommand{\ai}{\kw{ai}}
\newcommand{\no}{\kw{no.}}
\newcommand{\pid}{\kw{pid}}
\newcommand{\sid}{\kw{sid}}
\newcommand{\zip}{\kw{zip}}
\newcommand{\addr}{\kw{addr}}
\newcommand{\ph}{\kw{ph}}
\newcommand{\gndr}{\kw{gndr}}
\newcommand{\ppref}{\kw{ppref}}
\newcommand{\dob}{\kw{dob}}
\newcommand{\epref}{\kw{epref}}
\newcommand{\imei}{\kw{imei}}
\newcommand{\model}{\kw{model}}
\newcommand{\ipl}{\kw{ipl}}
\newcommand{\shno}{\kw{shno.}}
\newcommand{\dai}{\kw{dai}}
\newcommand{\shai}{\kw{shai}}
\newcommand{\dno}{\kw{dno.}}
\newcommand{\srch}{\kw{srch}}
\newcommand{\visit}{\kw{visit}}
\newcommand{\eid}{\kw{eid}}
\newcommand{\knull}{\kw{null}}

\newcommand{\seller}{\kw{seller}}
\newcommand{\type}{\kw{type}}
\newcommand{\fee}{\kw{fee}}
%\newcommand{\city}{\kw{city}}
\newcommand{\weight}{\kw{weight}}
\newcommand{\product}{\kw{product}}
\newcommand{\descr}{\kw{description}}
\newcommand{\quantity}{\kw{quantity}}
\newcommand{\price}{\kw{price}}
\newcommand{\online}{\kw{online\_year}}
\newcommand{\create}{\kw{date\_created}}
\newcommand{\sale}{\kw{on\_sale\_product}}
\newcommand{\shop}{\kw{shop}}
\newcommand{\items}{\kw{item}}
\newcommand{\delivery}{\kw{delivery}}

\newcommand{\CMap}{\kw{CMap}}
\newcommand{\CExtend}{\kw{CExtend}}
\newcommand{\Fix}{\kw{Fix}}
\newcommand{\Validate}{\kw{Validate}}

\newcommand{\enhance}{\kw{Enhance}}
\newcommand{\denhance}{\kw{DEnhance}}
\newcommand{\tid}{\kw{tid}}
%\newcommand{\incchase}{{\mbox{\small\sf IncChase}\xspace}\xspace}
%\newcommand{\pchase}{{\mbox{\small\sf PChase}\xspace}\xspace}

\newcommand{\incchase}{\kw{IncChase}}
\newcommand{\pchase}{\kw{PChase}}

\newcommand{\reviseW}[1]{#1}
\newcommand{\reviseWR}[1]{{\color{red}{#1}}}
\newcommand{\reviseR}[1]{#1}
\newcommand{\reviseB}[1]{#1}
%\newcommand{\cwarn}[1]{{\color{red}{#1}}}
\newcommand{\cwarn}[1]{#1}

\newcommand{\paper}{\kw{Paper}}
\newcommand{\finance}{\kw{Finance}}
\newcommand{\synthetic}{\kw{Syn}}
\newcommand{\movie}{\kw{Movie}}


\newcommand{\err}{\kw{err}}
\newcommand{\dup}{\kw{dup}}
\newcommand{\rlv}{\kw{rlv}}

\newcommand{\imdbm}{\kw{imdbm}}
\newcommand{\tmdbm}{\kw{tmdbm}}
\newcommand{\prdb}{\kw{prdb}}
\newcommand{\myear}{\kw{year}}
\newcommand{\budget}{\kw{budget}}
\newcommand{\orilan}{\kw{language}}
\newcommand{\dirct}{\kw{director}}
\newcommand{\mtitle}{\kw{title}}
\newcommand{\releasetime}{\kw{time}}
\newcommand{\length}{\kw{len}}
\newcommand{\mcountry}{\kw{country}}
\newcommand{\birthyear}{\kw{birth}}
\newcommand{\deathyear}{\kw{death}}
\newcommand{\profession}{\kw{job}}
\newcommand{\coworker}{\kw{coworker}}
\newcommand{\cast}{\kw{cast}}
\newcommand{\director}{\kw{director}}
\newcommand{\movies}{\kw{movies}}

\newcommand{\BigDansing}{\kw{BigDansing}}
\newcommand{\ProgressiveER}{\kw{ProgressiveER}}
\newcommand{\GGCP}{\kw{GGCP}}
\newcommand{\Sd}{{\mathcal S}}

\newcommand{\BCQE}{\kw{BCQE}}

\newcommand{\HTP}{\kw{HTP}}

\newcommand{\CGDs}{\kw{CGDs}}

%\newcommand{\ie}{\emph{i.e.,}\xspace}


\newcommand{\MRM}{\kw{MRLMatch}}
\newcommand{\SE}{\kw{SMatch}}
\newcommand{\PE}{\kw{PMatch}}
\newcommand{\IE}{\kw{IMatch}}

\newcommand{\Local}{\kw{LocalMatch}}
\newcommand{\TC}{\kw{Closure}}
\newcommand{\Mg}{\kw{Merge}}

\newcommand{\MHSP}{\kw{MHFP}}

\newcommand{\CQ}{\kw{CQ}}
\newcommand{\CQs}{\kw{CQs}}

\newcommand{\freq}{\kw{freq}}


\newcommand{\HQ}{\kw{HyperCube}}
\newcommand{\PT}{\kw{HyPart}}

\newcommand{\car}{\kw{Car}}

\newcommand{\HA}{\kw{HashAssign}}

\newcommand{\MN}{\kw{MergeNode}}

\newcommand{\tabincell}[2]{
\begin{tabular}{@{}#1@{}}#2\end{tabular}
}

%\floatsetup[table]{capposition=top}



\newcommand{\HC}{\kw{HC}}
\newcommand{\MQO}{\kw{MQO}}
\newcommand{\SeqM}{\kw{Match}}

\newcommand{\HGIN}{\kw{HGIN}}


\newcommand{\IDGNN}{\kw{IDGNN}}

\newcommand{\GNNs}{\kw{GNNs}}


\newcommand{\HFrame}{\kw{HFrame}}

%\newcommand{\C}{{\mathcal C}}

\newcommand{\cdd}[1]{|\!|{#1}|\!|}

\newcommand{\dual}{\kw{DualSim}}


\newcommand{\SHP}{\kw{SHP}}
\section{Introduction}
\label{sec-intro}
Graph pattern matching is 
an important problem in computer 
science, focusing on identifying instances of a pattern $Q$ 
within a graph $G$.
This topic has received 
extensive interest across various fields,
\eg 
discovering motifs in biological network
~\cite{chen2019hogmmnc, licheri2021grapes, tian2007saga}, 
finding experts in social networks
~\cite{shi2017multi, fei2013expfinder, brynielsson2010detecting}, 
and reasoning on knowledge graph
~\cite{xu2019cross, teru2020inductive, li2022hismatch}.

There exist two semantics for graph pattern matching, 
namely subgraph isomorphism and homomorphism.
(1) The isomorphic semantics demands that
the matched instance of $Q$  in 
$G$ \revise{(\ie a subgraph of $G$)} must be the same as $Q$.
The semantics has been widely studied, 
\eg discovering motifs from biological network~\cite{chen2019hogmmnc, licheri2021grapes, tian2007saga}.
(2) Homomorphic semantics,
on the other hand, permits that 
pattern $Q$ can be ``embedded'' into the 
identified \revise{subgraph of $G$}, 
allowing the  
\revise{subgraph} to be deviated from $Q$.
It has been widely used in graph query languages 
like Cypher~\cite{Cypher} and SPARQL~\cite{SPARQL}.
The difference between isomorphism and homomorphism
lies in the possibility 
of allowing different vertices in $Q$
can be mapped to the same vertex in $Q$ (see Section~\ref{sec-pre}).
\looseness = -1


\stitle{Related work}. We categorized the related work as follows.

\etitle{Graph pattern matching algorithm}.
Algorithms for subgraph isomorphism have been
extensively studied.
Ullmann proposed the first backtracking-based algorithm for subgraph isomorphism~\cite{Iso-match-6}. 
Subsequent research in algorithms for exact graph pattern 
matching focuses on pruning methods~\cite{carletti2017introducing,carletti2019vf3}, 
the application of index~\cite{CECI,sun2020subgraph}, 
index methods based on data mining~\cite{yan2004graph,zhao2007graph}, and so on. 
\looseness=-1

Isomorphic mappings suffers from two major drawbacks:
(1) the decision version of subgraph isomorphism is 
NP-complete~\cite{cook2023complexity},
(2) The requirement of  
\revise{injective} map 
could be too strict in many cases,
especially when the data source is noisy~\cite{yuan2011efficient}. 
 %
To address these challenges, 
(1) a series of studies focus on designing similarity functions and 
deriving matching results through  
top-k candidates or the establishment of a similarity threshold
\cite{fan2010graph,cheng2013top,dutta2017neighbor,wu2013ontology,khan2013nema,zhao2013partition}. 
(2) The other studies relax the  
\revise{injective} map to less restrictive  
relations.
Bounded simulation 
is developed to bound the lengths of 
matching paths~\cite{fan2010grapha}.
\revise{Dual simulation and } 
Strong simulation \revise{are proposed}
 in \cite{ma2011capturing}  to 
\revise{bound directions of edges 
	and 
the distance of vertices in a match, 
respectively}. 
Subsequently, various simulation-based
matching algorithms have been developed~\cite{fan2013incremental,fan2014distributed, fard2014distributed,gao2016prs,du2018personalized}.
\looseness = -1


\etitle{ML Models for graph pattern matching}.
Multiple ML models are developed for graph pattern matching~\cite{NeuroMatch,Sub-ML-1,Sub-ML-2,Sub-ML-3,Sub-ML-4}.
Most of these models are based on Graph Neural Networks (GNNs),
to compute embeddings of vertices in a graph~\cite{scarselli2008graph},
and make a prediction using the order-embedding space~\cite{NeuroMatch}
or MLP~\cite{Sub-ML-1}.
Most work focus on 
enhancing GNNs architecture~\cite{wang2022twin,wang2022equivariant,zhang2021eigen,bouritsas2022improving,feng2022powerful},
to learn more topological information. 
Zhang et al. \cite{zhang2023expressive} 
encoded distance information into vertex embeddings. 
Geerts et al. \cite{geerts2022expressiveness} 
proposed a model to aggregate information 
from $k$-order neighbors using computational operators. 
%
You et al. \cite{you2021identity} injected a unique 
color identity to nodes, 
and performed heterogeneous message passing 
to obtain the node embedding. 
Xu et al. \cite{xu2018powerful} 
optimized the aggregation function 
by using multisets, 
to generate distinct embeddings for different nodes. 
Maron et al. \cite{maron2018invariant} 
proposed a model consisting of 
both equivariant linear layers and nonlinear activation layers. 

There has also been work on
studying properties of GNNs.
(1) The expressive power of these GNN
models can be bounded by 
the Weisfeiler-Lehman (WL) hierarchy~\cite{xu2018powerful,morris2019weisfeiler,barcelo2022weisfeiler,zhang2023expressive,GNN-WL-C},
descriptive complexity~\cite{Express-LICS},
fragments of FO logic~\cite{Expressive-FOC2}
and global feature map transformer~\cite{Express-GFMT}. 
See~\cite{zhang2023expressive} for a recent survey.
(2) The generalization error
bound for GNNs can be measured by 
Rademacher complexity~\cite{Gen-bound}, VC dimension~\cite{VC-GNN}
and generalization gap~\cite{Generalization-gap,Generalization-gap}.

\section{Preliminaries}
\label{sec-pre}
Let $\Gamma$ be a countably 
infinite set of labels.

\stitle{Graphs}. 
Consider directed labeled graphs $G=(V, E, L)$,
where (1) $V$ is a finite set of vertices;
(2) $E\subseteq V\times \Gamma \times V $ is a finite set 
of edges, 
where $\langle v_1, r, v_2\rangle$ is an edge
from $v_1$ to $v_2$,  
\revise{carrying} label $r\in \Gamma$;
and (3) $L$ is a labeling function that 
assigns a label $L(v) \in \Gamma$
to each vertex $v$ in $V$.
Let $R_G$ be the set of all edge labels in $G$.
\looseness=-1

For each vertex $u {\in} V$, 
denote by $N^+_r(u)$ (resp. $N^-_r(u)$) the set of 
its  \revise{{\em outgoing} (resp. {\em incoming}) neighbors} $u_1$ of $u$
\revise{such that there exists} an  
edge labeled $r$ from $u$ to $u_1$
(resp.  
from $u_1$ to $u$).
That is, $N^+_r(u) {=} \{u_1\mid \langle u, r, u_1\rangle {\in} E\}$
and $N^-_r(u) {=} \{u_1\mid \langle u_1, r, u\rangle {\in} E\}$.
Denote by $N(u)$ the set of all neighbors of $u$, 
\ie $N(u) {=} \bigcup_{r\in R_G} (N^-_r(u) \cup N^+_r(u))$.

 

\stitle{Isomorphism and homomorphism.} 
We introduce the semantics of 
subgraph
isomorphism and homomorphism.
\revise{Consider a pattern  
$Q{=} (V_Q, E_Q, L_Q)$} and a graph $G {=} (V_G, E_G, L_G)$.
Here, a pattern is also a directed 
labeled graph, introduced
to simplify the presentation.  
\looseness=-1

 
\etitle{Subgraph isomorphism}.
An \textit{isomorphic mapping} from  
\revise{$Q$ to $G$},
denoted by 
\revise{$Q \sqsubseteq G$} is an injective mapping $\varphi$
from $V_Q$ to $V_G$  
such that
(1) for each $u{\in} V_Q$, $u$ and $\varphi(u)$ 
carry the same label, \ie $L_Q(u) = L_G(\varphi(u))$; and
(2) the adjacency relation \revise{preserves},  
\ie~\revise{$\langle u, r, u'\rangle \in E_Q$ if and only 
	if $\langle \varphi(u), r, \varphi(u')\rangle \in E_G$}. 
When such an injective $\varphi$ exists, 
we say that $Q$ is subgraph isomorphic to $G$. 
 
 
\section{Framework for homomorphic mapping}
\label{sec-model}


We provide \HFrame, a framework~\revise{for the subgraph homomorphic problem,}
by integrating algorithmic techniques and machine learning models.
The framework is outlined in Figure~\ref{fig:workflow}.
Given a graph $G$ and a pattern $Q$,
(1) it first leverages \kw{DualSim}
to refine the set ${\mathcal C}(u)$ of 
candidates for each vertex $u$ in $Q$ and subsequently constructs $S_{(Q, G)}$.
Recall that ${\mathcal C}(u)$ initially 
	consists of vertices in $G$ that carry the same label as $u$
	(see Section~\ref{sec-intro}).
%
(2) Then it checks whether ${\mathcal C}(u)$
is empty \revise{for some vertex $u$ in $Q$};
if so, then it returns \false, \ie 
no homomorphic mapping from $Q$ to $G$ exists;
otherwise,  
it constructs a subgraph $G'$ of $G$ 
	using only vertices appearing in $S_{(Q, G)}$; 
intuitively, it removes irrelevant vertices and edges from $G$,
which improves the performance of \HGIN (see Section~\ref{exp:sec-exp}).
(3) Finally, it  picks a {\em pivot vertex} $u_0$ from $Q$, and invokes ML model \HGIN
to check  
the existence of a homomorphic mapping  
$\varphi$ from $Q$ to $G'$
 such that $\varphi(u_0)=v$
for each 
vertex $v$ in 
${\mathcal C}(u_0)$, and 
	returns \true if \HGIN predicts 
	that such a mapping exists for some $v$ in $G'$.



  


\input{sec-proof}
%\section{Framework}
We propose a framework to detect homomorphism
by combining simulation and GNN.

give a figure for the framework

\stitle{Overview}.

\stitle{Properties}.

%\vspace{-2ex}
\section{Experimental study}
\label{exp:sec-exp}



\section{Conclusion}
In this paper, we proposed a framework \HFrame
for the subgraph homomorphic problem
by combining algorithmic solutions
and ML models.
We showed that \HFrame is more expressive
that the vanilla GNN models for sugbraph matching.
Moreover, we also provided the first generalization 
error
bound for the subgraph matching problem.
Using real-life and synthetic graphs, 
experimental results validate the efficiency
and effectiveness of the proposed frameworks.
\looseness = -1


%}

%%
%% The acknowledgments section is defined using the "acks" environment
%% (and NOT an unnumbered section). This ensures the proper
%% identification of the section in the article metadata, and the
%% consistent spelling of the heading.
%\clearpage
%\balance
%\input{reference}


%%
%% If your work has an appendix, this is the place to put it.
%\appendix


%\clearpage
%\balance
% Reference format: follow the numeric style used in
% "Automatic cable harness layout routing in a customizable 3D environment".
\bibliographystyle{elsarticle-num}
\bibliography{paper}



%\clearpage
%\balance
%\onecolumn
%\begin{verbatim}
	\appendix
	\input{appendix}
	%
%\subsection*{Proof of Theorem~\ref{thm-generalization-bound}}
\subsection*{A4. Proof of Theorem~\ref{thm-generalization-bound}}
We establish the Rademacher complexity  of \HFrame
by extending the proof of 
the Rademacher complexity for \GNNs~\cite{Gen-bound}.
Since \dual does not have any parameters, 
and works for any inputs, we only 
bound the Rademacher complexity of \HGIN.
We prove the bound in multiple steps:
(1) represent the computation of \HGIN
by a tree;
(2) quantify the impact of change 
of parameters on the embeddings;
(3) bound the changes in prediction probability;
(4) bound the covering numbers for \HGIN;
and (5) finally bound the Rademacher complexity
using the covering numbers.

In the following, we
show these steps one by one. 
In the proof of Theorem~\ref{thm-upper-bound}, 
the computation of \HGIN is represented by a 
computation tree.
It remains to prove steps (2)-(5).	
	
\stitle{$\bullet$ Step 2: Bound the impact of parameters}.
We first construct a formula 
to represent changes of embeddings, and then bound the changes.


Based on Equations (1)-(4), 
we can deduce the following upper bound
for the changes of embeddings:
\begin{small}
\begin{align*}
\Delta_{L}  \triangleq & \quad\|T_{L}(W_{1},\mathds{W})-T_{L}(W_{1}^{\prime},\mathds{W}^{\prime})\|_{2} \\
  =&\quad\bigg\|\kw{COM}(W_{1}x_{L} + \sum_{r\in R} W_r\underbrace{\rho(\sum_{j\in  N^+_r(x_{L})}\kw{MSG}_{{\bf 1}[j=v]}(T_{L-1,j}(W_{1},\mathds{W}))}_{\overset{\Delta}
 {\operatorname*{\operatorname*{=}}}R^+(W_{1},\mathds{W},x_{L},r)}),\\
 &\quad  W_{1}x_{L}+\sum_{r\in R} W_r\underbrace{\rho(\sum_{j\in 
 		N^-_r(x_{L})}\kw{MSG}_{{\bf 1}[j=v]}(T_{L-1,j}(W_{1},\mathds{W}))}_{\overset{\Delta}
 	{\operatorname*{\operatorname*{=}}}R^-(W_{1},\mathds{W},x_{L},r)})\\
 & -\quad\kw{COM}(W'_{1}x_{L}+\sum_{r\in R} W'_r\rho(\sum_{j\in 
 		N^+_r(x_{L})}\kw{MSG}_{{\bf 1}[j=v]}(T_{L-1,j}(W'_{1},\mathds{W}'))), \\
 		&\quad  W'_{1}x_{L}+\sum_{r\in R} W'_r\rho(\sum_{j\in 
 		N^-_r(x_{L})}\kw{MSG}_{{\bf 1}[j=v]}(T_{L-1,j}(W'_{1},\mathds{W}')))\bigg\|_{2} \\
  \leq& \quad 2C_{\kw{COM}}\left\|(W_{1}-W_{1}^{\prime})x_{L}\right\|_{2} \\
 & + C_{\kw{COM}}\sum_{r\in R}\|(W_rR^+(W_{1},\mathds{W},x_{L})-W_r^{\prime}R^+(W_{1}^{\prime},\mathds{W}^{\prime},x_{L}))\|_{2}\\
 & +\quad C_{\kw{COM}}\sum_{r\in R}\|(W_rR^-(W_{1},\mathds{W},x_{L})-W_r^{\prime}R^-(W_{1}^{\prime},\mathds{W}^{\prime},x_{L}))\|_{2}  \tag{10}\\
  \leq& \quad 2C_{\kw{COM}}\left\|(W_{1}-W_{1}^{\prime})x_{L}\right\|_{2} \\
& +\quad C_{\kw{COM}}\sum_{r\in R}\|(W_rR^+(W_{1},\mathds{W},x_{L})-W_r^{\prime}R^+(W_{1},\mathds{W},x_{L}))\|_{2}\\
& + C_{\kw{COM}}\sum_{r\in R}\|(W^{\prime}_rR^+(W_{1},\mathds{W},x_{L})-W_r^{\prime}R^+(W_{1}^{\prime},\mathds{W}^{\prime},x_{L}))\|_{2}\\
& +\quad C_{\kw{COM}}\sum_{r\in R}\|(W_rR^-(W_{1},\mathds{W},x_{L})-W_r^{\prime}R^-(W_{1},\mathds{W},x_{L}))\|_{2}\\
& + C_{\kw{COM}}\sum_{r\in R}\|(W^{\prime}_rR^-(W_{1},\mathds{W},x_{L})-W_r^{\prime}R^-(W_{1}^{\prime},\mathds{W}^{\prime},x_{L}))\|_{2} \tag{11}
\end{align*}
\end{small}

Equation~(10) holds, since we assume that 
 function \kw{COM} is $C_{\kw{COM}}$-Lipschitz.
 
 


To bound $\Delta_{L}$, 
we bound the following:
\begin{itemize}
\item[(a)] $\|(W_{1}{-}W_{1}^{\prime})x_{L}\|_{2}$;
\item[(b)]  $\|W_rR^+(W_{1},\mathds{W},x_{L}){-}W_r^{\prime}R^+(W_{1},\mathds{W},x_{L})\|_{2}$;
\item[(c)]  $\|W^{\prime}_rR^+(W_{1},\mathds{W},x_{L}){-}W_r^{\prime}R^+(W_{1}^{\prime},\mathds{W}^{\prime},x_{L})\|_{2}$;
\item[(d)]  $\|W_rR^-(W_{1},\mathds{W},x_{L})-W_r^{\prime}R^-(W_{1},\mathds{W},x_{L})\|_{2}$;
and 
\item[(e)]  $\|(W^{\prime}_rR^-(W_{1},\mathds{W},x_{L}){-}W_r^{\prime}R^-(W_{1}^{\prime},\mathds{W}^{\prime},x_{L})\|_{2}$.
\end{itemize}
%\looseness=-1

For term (a), we can assume an upper bound
for the norm of $W_1$ and $x_L$. 
For (b), 
we have 
$\|W_rR^+(W_{1},\mathds{W},x_{L})-W_r^{\prime}R^+(W_{1},\mathds{W},x_{L})\|_{2}
\leq \|(W_r-W_r^{\prime})R^+(W_{1},\mathds{W},x_{L})\|_{2}$.
Then we only need to bound $\|R^+(W_{1}^{\prime},\mathds{W}^{\prime},x_{L})\|_{2}$.
For (c), 
we have that 
$\|W^{\prime}_rR^+(W_{1},\mathds{W},x_{L})-W_r^{\prime}R^+(W_{1}^{\prime},\mathds{W}^{\prime},x_{L})\|_{2}
\leq \|W^{\prime}_r(R^+(W_{1},\mathds{W},x_{L})-R^+(W_{1}^{\prime},\mathds{W}^{\prime},x_{L}))\|_{2}$.
Therefore, we bound 
$\|R^+(W_{1},\mathds{W},x_{L})-R^+(W_{1}^{\prime},\mathds{W}^{\prime},x_{L})\|_{2}$.
The cases (d) and (e) can be analyzed similarly.
In the following, we only provide
upper bounds for $\|R^+(W_{1}^{\prime},\mathds{W}^{\prime},x_{L})\|_{2}$
and $\|R^+(W_{1},\mathds{W},x_{L})-R^+(W_{1}^{\prime},\mathds{W}^{\prime},x_{L})\|_{2}$.

\vspace{2ex}

\eat{
analyze the terms containing $R^+(W_{1},\mathds{W},x_{L})$.
The terms containing $R^+(W_{1},\mathds{W},x_{L})$ 
can be analyzed similarly.
}


\etitle{Upper bound for $\|R^+(W_{1},\mathds{W},x_{L})-R^+(W_{1}^{\prime},\mathds{W}^{\prime},x_{L}))\|_{2}$}.
We start with the upper bounds for term $\|R^+(W_{1},\mathds{W},x_{L})-R^+(W_{1}^{\prime},\mathds{W}^{\prime},x_{L}))\|_{2}$.
Based on the definition of 
$R^+(W_{1},\mathds{W},x_{L})$, we have the following.

\eat{
\begin{equation*}\begin{aligned}
&\|W_r^{\prime}R^+(W_1,\mathds{W},x_L)-W_r^{\prime}R^+(W_1^{\prime},\mathds{W}^{\prime},x_L)\|_2 
\leq&\|W_r^{\prime}\|_2\|R^+(W_1,\mathds{W},x_L)-R^+(W_1^{\prime},\mathds{W}^{\prime},x_L)\|_2
\end{aligned}\end{equation*}
}


\begin{scriptsize}	
%\vspace{-7ex}
\begin{align*}
&\|R^+(W_1,\mathds{W},x_L,r)-R^+(W_1',\mathds{W}',x_L,r)\|_2\\ % \tag{eq-1} 
  &\leq C_\rho\bigg\|\sum_{j\in N_r^+(x_L)}\bigg(\kw{MSG}_{{\bf 1}[j=v]}(T_{L-1,j}(W_1,
  \mathds{W})) -\kw{MSG}_{{\bf 1}[j=v]}(T_{L-1,j}(W_1',\mathds{W}'))\bigg)\bigg\|_2  \tag{12}
  \\
 & \leq C_\rho\sum_{j\in N_r^+(x_L)}\bigg\|\bigg(\kw{MSG}_{{\bf 1}[j=v]}(T_{L-1,j}(W_1,
  \mathds{W})) -\kw{MSG}_{{\bf 1}[j=v]}(T_{L-1,j}(W_1',\mathds{W}'))\bigg)\bigg\|_2 \nonumber
  \\
  &\leq C_\rho C_{\kw{MSG}_0}C_{\kw{MSG}_1}\sum_{j\in N_r^+(x_L)}\bigg\|T_{L-1,j}(W_1,\mathds{W})-
  T_{L-1,j}(W_1',\mathds{W}')\bigg\|_2 \tag{13}\\
  &= C_\rho C_{\kw{MSG}_0}C_{\kw{MSG}_1}\sum_{j\in N_r^+(x_L)}\Delta_{L-1,j}\\
  &\leq C_\rho C_{\kw{MSG}_0}C_{\kw{MSG}_1}d\max_{j\in N_r^+(x_L)}\Delta_{L-1,j}\tag{14}
\end{align*}
\end{scriptsize}


Equations (12) and (13) follow from the assumption 
that Functions $\rho$, $\kw{MSG}^{k,r}_0$ and 
$\kw{MSG}^{k,r}_0$ are $C_\rho$-Lipschitz, 
$C_{\kw{MSG}_0}$-Lipschitz  and $C_{\kw{MSG}_1}$-Lipschitz,
respectively.
Equation (14) holds, since we assume that 
the degrees of vertices in $G$ are bounded by 
a constant $d$.
\looseness = -1

\eat{
Since we assume that $\|W'_r\|\leq B_2$,
we can deduce that the bound as follows.

\begin{equation*}\begin{aligned}
||W_r^{\prime}R(W_1,\mathds{W},x_L,r)  - W_r^{\prime}R(W_1^{\prime},\mathds{W}^{\prime},x_L,r)\|_2
 &\leq B_2C_\rho C_{\kw{MSG}_0}C_{\kw{MSG}_1}\sum_{j\in N_r^+(x_L)}\Delta_{L-1,j} \\
 &\leq B_2C_\rho C_{\kw{MSG}_0}C_{\kw{MSG}_1}d\max_{j\in N_r^+(x_L)}\Delta_{L-1,j}
\end{aligned}\end{equation*}


\begin{equation*}\begin{aligned}
\Delta_{L} & \leq\quad 2C_{\kw{COM}} B_x\left\|(W_1-W_1^{\prime})\right\|_2 
  +\quad C_{\kw{COM}} |R| B_2C_\rho C_{\kw{MSG}_0}C_{\kw{MSG}_1}d\max_{j\in N_r^+(x_L)}\Delta_{L-1,j} \\
 &\quad +\kw{COM}|R|\max_{r\in R}\|(W_r-W_r^{\prime})R^+(W_1,\mathds{W},x_L,r)||_2
  + \kw{COM}|R|\max_{r\in R}\|(W_r-W_r^{\prime})R^-(W_1,\mathds{W},x_L,r)||_2
\end{aligned}\end{equation*}
}

Observe that the upper bound for
$\Delta_{L}$ depends on that of $\|R^+(W_1,\mathds{W},x_L,r)-R^+(W_1',\mathds{W}',x_L,r)\|_2$,
which in turn dependents on $\Delta_{L-1,j}$.
Then we deduce a bound for $\|R^+(W_1,\mathds{W},x_L,r)-R^+(W_1',\mathds{W}',x_L,r)\|_2$
by expanding the bounds (see below). 


% \begin{equation}
% \begin{aligned}
% \|R(W_1,\mathds{W},x_L)\|_2 \\ 
% \leq\quad C_\rho C_gd\min\left\{b\sqrt{r},C_\phi B_1B_x\frac{(\mathcal{C}d)^L-1}{\mathcal{C}d-1}\right\}
% \end{aligned}
% \end{equation}



\etitle{Upper bounds for $\|R^+(W_{1}^{\prime},\mathds{W}^{\prime},x_{L}))\|_{2}$}.
Based on the definition of $\|R^+(W_{1},\mathds{W},x_{L},r)\|_{2}$, 
we can deduce the following.

\begin{align*}
  \|R^+(W_{1},\mathds{W},x_{L},r)\|_{2} 
 & =\bigg\|\rho\big(\sum_{j\in N_r^+(x_{L})}\kw{MSG}_{{\bf 1}[j=v]}(T_{L-1,j}(W_{1},\mathds{W})
 \big)\big)\bigg\|_{2} \\
 & \leq C_\rho\bigg\|\sum_{j\in N_r^+(x_L)}\kw{MSG}_{{\bf 1}[j=v]}(T_{L-1,j}(W_1,\mathds{W}))\bigg\|
 _2 \\
 & \leq C_\rho\sum_{j\in N_r^+(x_L)}\bigg\|\kw{MSG}_{{\bf 1}[j=v]}(T_{L-1,j}(W_1,
 \mathds{W}))\\
 &\quad\quad -\kw{MSG}_{{\bf 1}[j=v]}(0)\bigg\|_2 \\
 & \leq C_{\rho}C_{\kw{MSG}_0}C_{\kw{MSG}_1}\sum_{j\in N_r^+(x_{L})}\bigg\|T_{L-1,j}(W_{1},\mathds{W})
 \bigg\|_{2} \\
 & \leq C_\rho C_{\kw{MSG}_0}C_{\kw{MSG}_1}d\max_{j\in N_r^+(x_L)}\bigg\|T_{L-1,j}(W_1,\mathds{W})\bigg\|_2
\end{align*}


The equations hold due to the reasons given above.
Observe that $\|R^+(W_{1},\mathds{W},x_{L},r)\|_{2}$ 
is bounded using $\bigg\|T_{L-1,j}(W_1,\mathds{W})\bigg\|_2$,
\ie the changes of the predictions of neighbors.
We continue to expand $\bigg\|T_{L-1,j}(W_1,\mathds{W})\bigg\|_2$
as follows.


\begin{align*}
		&\quad\quad\bigg\|T_{L-1,j}(W_1,\mathds{W})\bigg\|_2\\
		 &=\quad\bigg\|\kw{COM}(W_{1}x_{L}+\sum_{r\in R} W_r R^+(W_1,W_2,x_{L-1,j},r),
		\\
		&\quad\quad\quad W_{1}x_{L}+\sum_{r\in R} W_r R^-(W_1,W_2,x_{L-1,j},r)) \bigg\|
_2 \\
		& =\quad\bigg\|\kw{COM}(W_{1}x_{L}+\sum_{r\in R} W_r R^+(W_1,W_2,x_{L-1,j},r),
W_{1}x_{L}\\
&\quad\quad\quad+\sum_{r\in R} W_r R^-(W_1,W_2,x_{L-1,j},r)) -
\kw{COM}(0,0) \bigg\|
_2 \\
		& \leq\quad C_{\kw{COM}}\bigg\|W_{1}x_{L-1,j}+\sum_{r\in R} W_r 
		R^+(W_{1},W_{2},x_{L-1,j},r)\bigg\|_{2} \\
		&\quad\quad+  C_{\kw{COM}}\bigg\|W_{1}x_{L-1,j}+\sum_{r\in R} W_r 
R^-(W_{1},W_{2},x_{L-1,j},r)\bigg\|_{2} \tag{15} \\
		& \leq\quad 2C_{\kw{COM}}\bigg\|W_{1}x_{L-1,j}\bigg\|_{2}+C_{\kw{COM}}\sum_{r\in R} \bigg\|
		W_rR^+(W_{1},W_{2},x_{L-1,j},r)\bigg\|_{2} \\
		&\quad\quad+ C_{\kw{COM}}\sum_{r\in R} \bigg\|
		W_rR^-(W_{1},W_{2},x_{L-1,j},r)\bigg\|_{2}\\
		& \leq\quad 2C_{\kw{COM}}B_{1}B_{x}+C_{\kw{COM}}|R|B_{2}\bigg\|
		R^+(W_{1},W_{2},x_{L-1,j},r)\bigg\|_{2}\\
		&\quad\quad+C_{\kw{COM}}|R|B_{2}\bigg\|
		R^-(W_{1},W_{2},x_{L-1,j},r)\bigg\|_{2}.  \tag{16}
\end{align*}


Equation (15) holds, since $\kw{COM}$
is $C_{\kw{COM}}$-Lipschitz.
And Equation (16) holds, since 
we assume that 
the norms of $W_1$ and $X_i$
is bounded by $B_1$ and $B_x$, respectively.
%
Using the above bound, we bound $\|R^+(W_1,\mathds{W},x_L,r)\|_2$ %can be bounded
as follows.

\begin{align*}
&  \|R^+(W_1,\mathds{W},x_L,r)\|_2 \\
 & \leq 2C_\rho C_{\kw{MSG}_0}C_{\kw{MSG}_1}C_{\kw{COM}} B_1B_xd  \\
 & +  C_\rho C_{\kw{MSG}_0}C_{\kw{MSG}_1}C_{\kw{COM}} |R| 
 B_2d\max_{j\in C(x_L)}\bigg(\bigg\|R^+(W_1,
 \mathds{W},x_{L-1,j})\bigg\|_2\\
 &\quad\quad+\bigg\|R^-(W_1,
 \mathds{W},x_{L-1,j})\bigg\|_2\bigg) \\
 & \leq 2C_\rho C_{\kw{MSG}_0}C_{\kw{MSG}_1}C_\phi B_1B_xd\sum_{\ell=0}^{L-1}(2C_\rho C_{\kw{MSG}_0}C_{\kw{MSG}_1}
 C_{\kw{COM}} B_2 |R| d)^\ell \\
 & =\quad 2C_\rho C_{\kw{MSG}_0}C_{\kw{MSG}_1}C_{\kw{COM}} B_1B_xd\frac{(\mathcal{C}d)^L-1}{\mathcal{C}d-1}. \tag{17}
\end{align*}

Let ${\mathcal C} = 2C_\rho C_{\kw{MSG}_0}C_{\kw{MSG}_1}C_{\kw{COM}} B_2 |R| d$.
Denote $M$ as $\frac{\left(\mathcal{C}d\right)^{L}-1}{\mathcal{C}d-1}$.
That is, $\|R^+(W_1,\mathds{W},x_L,r)\|_2  \leq {\mathcal C} B_1  B_xM$.

\eat{
Moreover, since we assume that $\|x\|_\infty=|x_1|+\ldots+|x_r|\leq 
b$ for each vector $x=(x_1, \ldots, x_r)\in R^r$, 
$\|x\|_2\leq \sqrt{x^2_1+\ldots + x^2_r} \leq \sqrt{r}\|x\|_\infty$
by the mean value inequality.
We have that 
$$\left\|T_{L-1,j}(W_1,\mathds{W})\right\|_2\leq b\sqrt{r}$$

Moreover, we have that 

\begin{align*}\|R^+(W_1,\mathds{W},x_L,r)\|_2\leq bdC_\rho C_{\kw{MSG}_0}C_{\kw{MSG}_1}\sqrt{r}
\end{align*}


Combining \tbf and \tbf, we have that
\begin{align*}\|R^+(W_1,\mathds{W},x_L,r)\|_2\leq C_\rho C_{\kw{MSG}_0}C_{\kw{MSG}_1}d\min\left\{b\sqrt{r},C_\kw{COM} 
B_1B_x\frac{(\mathcal{C}d)^L-1}{\mathcal{C}d-1}\right\}\end{align*}
}


\etitle{Upper bounds for $\Delta_{L}$}.
Combining Equations (11), (14) and (17), 
we get

\eat{
\begin{gather*}
\Delta_{L}\quad\leq\quad MB_{x}\left\|W_{1}-W_{1}^{\prime}\right\|_{2} 
+ M||R(W_1,\mathds{W},x_L)||_2||\mathds{W}-\mathds{W}^{\prime}||_2.
\end{gather*}
}

\begin{align*}
	\Delta_{L} & \leq\quad C_{\kw{COM}} B_x\|W_1-W_1^{\prime}\|_2 \\
	&\quad + 2C_{\kw{COM}} |R| B_2C_\rho C_{\kw{MSG}_0}C_{\kw{MSG}_1}d\max_{j\in C(x_L)}\Delta_{L-1,j} \\
	&\quad + C_{\kw{COM}}|R|\max_{r\in R}\|W_r-W_r^{\prime}\|_2\|R^+(W_1,\mathds{W},x_L,r)\|_2\\
	&\quad + C_{\kw{COM}}|R|\max_{r\in R}\|W_r-W_r^{\prime}||_2\|R^-(W_1,\mathds{W},x_L,r)\|_2\\
% & \leq\quad C_{\kw{COM}} B_x\|W_1-W_1^{\prime}\|_2 \\
%&\quad + 2C_{\kw{COM}} |R| B_2C_\rho C_{\kw{MSG}_0}C_{\kw{MSG}_1}d\max_{j\in C(x_L)}\Delta_{L-1,j} \\
%&\quad + 2 C_{\kw{COM}} |R| {\mathcal C} B_1  B_xM \max_{r\in R}\|W_r-W_r^{\prime}||_2\\
	&\leq\quad 4M \bigg(2B_{x}||W_{1}-W_{1}^{\prime}||_{2}\\
	&\quad+\max_{r\in R}\|W_r-
	W_r^{\prime}\|_2(\|R^+(W_1,\mathds{W},x_L,r)\|_2\\
	&\quad+\|R^-(W_1,\mathds{W},x_L,r)\|_2)\bigg) \tag{18}
\end{align*}



\eat{
\begin{equation}\Delta_L\leq MB_x\left|\left|W_1-W_1^{\prime}\right|\right|
_2+M\overline{R}||
W_r-W_r^{\prime}||_2\end{equation}
}

\stitle{$\bullet$ Step 3: Bound the prediction changes}.
We first reform the change of 
the prediction function as follows. 

\begin{equation*}\begin{aligned}
		P_{L} & \triangleq\quad\|T^Q_{L}(W_{1},\mathds{W})-T^G_{L}(W_{1},
		\mathds{W})\|_c \\
		& =\quad\bigg\|\kw{COM}(W_{1}x^Q_{L}+\sum_{r\in R} W_rR^+(W_{1},\mathds{W},x^Q_{L},r)), W_{1}x^Q_{L}\\
		& \quad\quad+\sum_{r\in R} W_rR^-(W_{1},\mathds{W},x^Q_{L},r))\\
			& \quad\quad-\kw{COM}(W_{1}x^G_{L}+\sum_{r\in R} W_rR^+(W_{1},\mathds{W},x^G_{L},r)), W_{1}x^G_{L}\\
			&\quad\quad+\sum_{r\in R} W_rR^-(W_{1},\mathds{W},x^G_{L},r)) \bigg\|_c\\
%		&=W_{1}(x^Q_{L}-x^G_{L}) + \sum_{r\in R} W_r(R(W_{1},
%		\mathds{W},x^Q_{L},r)-R(W_{1},\mathds{W},x^G_{L},r))\bigg\|_c\\
\end{aligned}\end{equation*}

Then the change of prediction can be bounded as follows.

\begin{scriptsize}
\begin{align*}
		\nabla_{L} & \triangleq \quad |P_L - P'_L|\\
		& =\quad\bigg|\|T^Q_{L}(W_{1},\mathds{W})-T^G_{L}(W_{1},
		\mathds{W})\|_c - \|T^Q_{L}(W'_{1},\mathds{W}')-T^G_{L}(W'_{1},
		\mathds{W}')\|_c\bigg| \\
		& \leq \quad\bigg|\|T^Q_{L}(W_{1},\mathds{W})-T^Q_{L}(W'_{1},
		\mathds{W}')-  (T^G_{L}(W_{1},
		\mathds{W})-T^G_{L}(W'_{1},
\mathds{W}'))\|_c\bigg| \tag{19}\\
		& \leq \quad\|T^Q_{L}(W_{1},\mathds{W})-T^Q_{L}(W'_{1},
\mathds{W}')\|_2 +
\|T^G_{L}(W_{1},
\mathds{W})-T^G_{L}(W'_{1},
\mathds{W}')\|_2 \\
&\leq \quad 4M \bigg(2B_{x}||W_{1}-W_{1}^{\prime}||_{2}+M\max_{r\in R}\|W_r-
W_r^{\prime}\|_2(\|R^+(W_1,\mathds{W},x_L,r)\|_2\\
&\quad\quad+\|R^-(W_1,\mathds{W},x_L,r)\|_2)\bigg)\tag{20}
\end{align*}
\end{scriptsize}


Equation (19) holds, due to the triangle inequality 
of norm $\|\cdot\|_c$.
And Equation (20) holds due to Equation (18) above.

\eat{
Here, we show that $\|a\|_c - \|b\|_c\leq \|a-b\|_c$.
Assume that $\|a\|_c = a_i$ and $\|b\|_c=b_j$.
Then $b_i\leq b_j$ and $-b_i\geq -b_j$. 
Hence, $a_i-b_j\leq a_i-b_i \leq \max_{k\in [1, d]} a_i-b_i$.
That is, $\|a\|_c - \|b\|_c\leq \|a-b\|_c$.
}

\stitle{$\bullet$ Step 4: Bound the covering numbers}.
To bound the change in $\nabla_{L}$ by 
a given threshold $\epsilon$, we use the covering number.
Intuitively, given a set $\M$ of matrices
that have a bound for the norm $\|\cdot\|_F$,
\ie $\|M\|_F\leq \lambda$ for all $M\in \M$,
and a number $\eta$, its covering 
number ${\mathcal N}(\M, \eta, \|\cdot\|_F)$ is defined as the minimum number
of balls of radius $\eta$, needed to cover
all matrices in $\M$~\cite{Covering-number}.
To bound the covering number, 
we use a lemma from~\cite{Covering-1}, 
which states that 
$${\mathcal N}(\M, \eta, \|\cdot\|_F)\leq \bigg(1+\frac{2d\lambda}{\epsilon}\bigg)^{d^2}$$

Here, $\varepsilon$ is a given threshold, $\lambda$
is the bound for the norm of matrices in $\M$, 
and $d$ is the dimension 
of the matrix.

Since the bound in Equation (20) has two kinds
of learned matrices, \ie $W_1$ and matrices in $\mathds{W}$.
Using this lemma, we deduce that 


\begin{align*}\mathcal{N}\left(W_{1},\frac{\epsilon}{8MB_{x}},\|\cdot\|_{F}\right)
\leq\left(1+\frac{16MB_{x}B_{1}\sqrt{d}}{\epsilon}\right)^{d^{2}}\end{align*}

\begin{align*}\mathcal{N}\left(W_r,\frac{\epsilon}{8M{\mathcal R}},
	\|\cdot\|_{F}\right)\leq\left(1+\frac{16M{\mathcal R}B_{2}\sqrt{d}}{\epsilon}
	\right)^{d^{2}} \quad (r\in R)
\end{align*}

Here, ${\mathcal R} = \max(\|R^+(W_1,\mathds{W},x_L,r)\|_2,
\|R^-(W_1,\mathds{W},x_L,r)\|_2) \leq {\mathcal C} B_1  B_xM$.
Therefore, the overall covering number for \HGIN is bounded
by the production of all these covering numbers~\cite{Gen-bound}.
\begin{align*}P=\left(1+\frac{16M\sqrt d\max\{B_xB_1,{\mathcal R}B_2\}}
{\epsilon}\right)^{2|R|d^2}\end{align*}

Because $\varepsilon$ is sufficiently small, 
\eg $\epsilon<4M\sqrt{d}\max\{B_xB_1,{\mathcal R}B_2\}$,
we have that $\log P\leq 2|R|d^2\log\left(\frac{32M\sqrt{d}\max\{B_xB_1,{\mathcal R}
B_2\}}{\epsilon}\right)$.
We will use this upper bound to prove the Rademacher complexity.

\eat{
\begin{equation*}
	3r^2\log\left(1+\frac{6M\sqrt{r}\max\{B_xB_1,\overline{R}B_2\}}{\epsilon}\right)
\end{equation*}
\begin{equation*}
	\epsilon<4 M\sqrt{r}\max\{B_xB_1,\overline{R}B_2\}
\end{equation*}
\begin{equation*}
	3r^2\log\left(\frac{8M\sqrt r\max\{B_xB_1,\overline{R}B_2\}}{\epsilon}\right)
\end{equation*}
}

\stitle{$\bullet$ Step 5: Bound for the Rademacher complexity}.
We finally can bound the Rademacher complexity.
To this end, we also apply a lemma from~\cite{RB-1},
which states that 
given a function ${\mathcal F}$ taking values
in $[0,1]$, 
its Rademacher complexity 
$\hat{\mathcal{R}}$ 
is bounded by $\inf_{\alpha>0}\left(\frac{4\alpha}{\sqrt{N_T}}+\frac{12}{N_T}\int_{\alpha}^{\sqrt{N_T}}\sqrt{\log\mathcal{N}(\mathcal{F},\epsilon,\|\cdot\|_F)}d\epsilon\right)$,
where $\mathcal{N}(\mathcal{F},\epsilon,\|\cdot\|_F)$ is the covering 
number of $\mathcal{F}$, and
$N_T$ is the number of training data.


From Step-4, we know that the covering number ${\mathcal N}$
is bounded by $\left(1+\frac{16M\sqrt d\max\{B_xB_1,{\mathcal R}B_2\}}
{\epsilon}\right)^{2|R|d^2}$.

Therefore, we can bound the covering number as follows.

\begin{align*}
\log \mathcal{N}(\mathcal{F},\epsilon,\|\cdot\|_F)&\leq 
\log \mathcal{N}(P,\frac{\epsilon}{2},\|\cdot\|_F)\\
&
\leq 2|R|d^2\log\left(\frac{32M\sqrt{d}\max\{B_xB_1,\overline{R}
	B_2\}}{\epsilon}\right)
\end{align*}

Setting $\alpha  = \frac{2}{\sqrt{N_T}}$, we have that 
$\hat{\mathcal{R}}(\mathcal{T})\leq\frac{8}{N_T}+
\frac{12}{\sqrt{N_T}}\sqrt{2|R|d^2\log T}$,
where $T=16M\sqrt{{N_T}d}\max\{B_xB_1,{\mathcal R}B_2\}$,
and $\mathcal{I}$ is the class of functions
mapping the computation of embedding to the prediction function $\nabla_{L}$.

We can verify that the max margin loss function 
is 1-Lipschitz.
Then ${\hat{\mathcal{R}}(\mathcal{J}_\gamma)}
%\leq \hat{\mathcal{R}}
%_T(\mathcal{I})
\leq\frac{8}{N_T}+\frac{12d}{\sqrt{N_T}}
\sqrt{2|R|d^2|R|\log T}$.

\stitle{The 1-Lipschitz of the loss function}.
We next show that the max margin loss function 
is 1-Lipschitz, 
\ie $|{\mathcal L}(h_u,h_v)-{\mathcal L}(h'_u,h'_v)|\leq 
|{\mathcal M}(h_u,h_v)-{\mathcal M}(h'_u,h'_v)|$.
Since the loss functions for positive data and negative data
are different,
we consider all possible four cases:
(1) both $(h_u,h_v)$ and $(h'_u,h'_v)$ are positive;
(2) both $(h_u,h_v)$ and $(h'_u,h'_v)$ are negative;
(3) $(h_u,h_v)$ is positive and $(h'_u,h'_v)$ is negative; and
(4) $(h_u,h_v)$ is negative and $(h'_u,h'_v)$ are positive.



\stab
(1) If both $(h_u,h_v)$ and $(h'_u,h'_v)$ are positive, 
then their loss are $-\alpha$.
Hence, 
$|{\mathcal L}(h_u,h_v)-{\mathcal L}(h'_u,h'_v)| = 0\leq 
|{\mathcal M}(h_u,h_v)-{\mathcal M}(h'_u,h'_v)|$.

\eat{
\stab
(2) If both $(Q_1,G_1)$ and $(Q_2,G_2)$ are negative,
and the margins between $(Q_1,G_1)$ and $(Q_2,G_2)$ are smaller than
$\alpha$, then the inequality holds.

\stab
(3) If both $(Q_1,G_1)$ and $(Q_2,G_2)$ are negative,
and the margins between $(Q_1,G_1)$ and $(Q_2,G_2)$ are larger than 
$\alpha$, then the loss is 0, and the inequality holds.
}

\stab
(2) If both $(h_u,h_v)$ and $(h'_u,h'_v)$ are negative,
when both ${\mathcal M}(h_u,h_v)$
and ${\mathcal M}(h'_u,h'_v)$ are not larger than 
$\alpha$, we have that 

\begin{align*}
&{\mathcal L}(h_u,h_v)-{\mathcal L}(h'_u,h'_v)|\\
&=(\alpha -\|\max(0,h_u-h_v)\|^2_2) - (\alpha - \|\max(0,h'_u-h'_v)\|^2_2)\\
&=-\|h_u-h_v\|^2_2 + \|h'_u-h'_v\|^2_2 \\
&\leq |{\mathcal M}(h_u,h_v)-{\mathcal M}(h'_u,h'_v)| \tag{21}
\end{align*}

Equation (21) holds, since 
$\mathcal{M}(h_u, h_v) = \cdd{\max(0, h_u-h_v)}_2^2$.

\vspace{2ex}
When ${\mathcal M}(h_u,h_v)\geq \alpha$ or
 ${\mathcal M}(h'_u,h'_v)\geq  \alpha$,
 the proof is similar.











\stab
(3) If $(h_u,h_v)$ is positive and $(h'_u,h'_v)$ is negative,
and ${\mathcal M}(h'_u,h'_v)\leq \alpha$, we have the following

$\|\max(0,h_u - h_v)\|_2 - \alpha  - (\alpha - \|\max(0,h'_u - h'_v)\|)\\
=   \|h'_u - h'_v \|_2- 2\alpha\\
\leq \|h_u - h_v\|_2 +  \|h'_u - h'_v\|_2- 2\|h'_u - h'_v \|_2\\
\leq \|h_u - h_v\|_2 - \|h'_u - h'_v \|_2\\
\leq |{\mathcal M}(h_u,h_v)-{\mathcal M}(h'_u,h'_v)|$





\stab
(4) If $(h_u,h_v)$ is positive and $(h'_u,h'_v)$ is negative,
and ${\mathcal M}(h'_u,h'_v)\geq \alpha$, then 
we have 

$
\|\max(0,h_u - h_v)\|_2 - \alpha  - \max(0,\alpha - \|\max(0,h'_u - h'_v)\|)\\
= - \alpha\\
\leq |{\mathcal M}(h_u,h_v)-{\mathcal M}(h'_u,h'_v)|$

\subsection*{A5. Complex queries}

\begin{figure}[tb!]
	%	\vspace{-6.7ex}
	\centerline{\includegraphics[width=0.8\columnwidth]{fig/query-example-2.eps}}
	\centering
	\vspace{-1.4ex}
	\caption{Example of complex query}
	\label{fig:query}
	\vspace{-1.4ex}
\end{figure}

When patterns are complex (\eg with multiple cycles
and large number of vertices),
the exact algorithms can be much slow 
and \HFrame is significantly faster than 
them. 
Consider a complex query extracted from Citeseerx,
which is depicted in 
Figure~\ref{fig:query}.
\kw{SIM}-\kw{TD}, \kw{EJ} and \kw{PJ} take
2.691s, 11.164s and 13.472s respectively to conduct the query,
while \HFrame takes only 0.081s.

\reviserdy{
\subsection*{A6. The Application of \HFrame}
We focus on the decision problem of the subgraph homomorphism problem, 
which can be applied in frequent pattern mining, graph pattern depenedency discovery, 
graph query evaluations, and accelerating subgraph enumerations. 

We can extend our algorithms to generate exact subgraph instances in the following two ways:
(1) plugging in the enumeration algorithm: 
we can apply \HFrame to verify the candidates of the first pattern vertex in the matching order, 
and apply exact algorithm to enumerate on the vertices, on which \HFrame returns true;
or (2) applying \HFrame to filter candidates: 
we can apply \HFrame to check whether there exists a match for each vertex v in C(u) for each pattern vertex u, 
remove all vertices in C(u) that does not have a match, 
and finally run the exact enumeration algorithm on the subgraph deduced by the remained vertices. 
Taking PJ as an example, we adopt the second method, 
and reports the results on Citation and DBLP as \ref{tab:app-filter}. 
We can see that \HFrame trades off a marginal loss in accuracy for a significant gain in computational efficiency.
\begin{table}[tb]
	\caption{Effectiveness of \HFrame in candidate filtering}
	\label{tab:app-filter}
	\vspace{-1.5ex}
	\centering
	\resizebox{1.01\columnwidth}{!}{
		\begin{tabular}{c||c|c|c|c} \hline
			& \multicolumn{2}{c|}{\bf{Accuracy}}
                & \multicolumn{2}{c}{\bf{Efficiency (s)}}\\
			\hline
			\diagbox{\bf{Method}}{\bf{Dataset}} &Citation&DBLP&Citation&DBLP\\\hline
			\kw{Only PJ}              & 1.000 & 1.000 & 1.545 & 1.803 \\ %\hline
			\kw{\HFrame+PJ}              & 0.815 & 0.832 & 0.562 & 0.547 \\ \hline
		\end{tabular}
	}
	\vspace{-2ex}
\end{table}

Additionally, in frequent pattern mining and graph pattern dependency discovery, 
the quality of patterns $Q$ is usually defined by the number of vertices matching the pivot of $Q$, 
rather than the number of matches. 
We have added an application of \HFrame in frequent pattern mining algorithm SPMiner\cite{ying2024representation}. 
We evaluated the performance on Enzyme dataset same as SPMiner, 
and calculated the median anchored frequency and grouped runtime by pattern size in \ref{tab:app-mining}. 
We found that when plugging \HFrame in the existing frequent pattern mining, 
\HFrame can discover high-quality patterns in a reasonable time.

\begin{table}[b]
        \caption{Effectiveness of \HFrame in frequent pattern mining}
            \label{tab:app-mining}
        \vspace{-1.5ex}
        \centering
        \resizebox{1.01\columnwidth}{!}{
            \begin{tabular}{c|c|c|c|c|c|c} \hline
                    \multicolumn{2}{c|}{} & \multicolumn{5}{c}{Pattern size} \\ \cline{3-7}
                    \multicolumn{2}{c|}{} & 5 & 6 & 7 & 8 & 9 \\ \hline
                    \multirow{2}{*}{Support} & SPMiner+\HFrame & 14,688 & 14,445 & 13,571 & 12,109 & 9,723 \\
                             & SPMiner & 13,005 & 14,572 & 13,691 & 9,831 & 7,928 \\ \hline
                    \multirow{2}{*}{Runtime(s)} & SPMiner+\HFrame & 160.85 & 193.77 & 226.01 & 256.43 & 292.16 \\
                             & SPMiner & 151.89 & 182.38 & 212.66 & 244.02 & 275.40 \\ \hline
            \end{tabular}
        }
        \vspace{-2ex}
\end{table}
}
	
\eat{

\clearpage
\subsection*{A5. Complex queries}

\begin{figure}[tb!]
	%	\vspace{-6.7ex}
	\centerline{\includegraphics[width=0.8\columnwidth]{fig/query-example-2.eps}}
	\centering
	\vspace{-1.4ex}
	\caption{Example of complex query}
	\label{fig:query}
	\vspace{-1.4ex}
\end{figure}

There exist some 

\warn{Figure~\ref{fig:query} illustrates a complex query example 
	that poses significant challenges for exact algorithms.}
}



\eat{
\begin{align*}
		\hat{\mathcal{R}}_{\mathcal{T}}(\mathcal{J}_{\gamma})\leq\frac{4}{\gamma m}+
		\frac{16r}{\gamma\sqrt{m}}\sqrt{3\operatorname{log}Q},where \\
		Q=16\sqrt{m}M\sqrt{r}\max\{B_{x}B_{1},\overline{R}B_{2}\}, \\
		M=C_\phi\frac{\left(\mathcal{C}d\right)^L-1}{\mathcal{C}d-1}, \\
		\overline{R}\leq C_{\rho}C_{g_0}C_{g_1}d\operatorname*{min}\left\{b\sqrt{r},B_{1}B_{x}
		M\right\}.
\end{align*}

\begin{align*}e\leq 1\end{align*}
\warn{Since we normalize the embeddings}

Hence, the empirical Rademacher complexity of $\mathcal{J}_\gamma$ with respect to $\mathcal{T}$ is
\begin{align*}\hat{\mathcal{R}}_{\mathcal{T}}(\mathcal{I})\leq\inf_{\alpha>0}\left(\frac{4\alpha}{\sqrt{m}}+\frac{12}{m}\int_{\alpha}^{2e\sqrt{m}}\sqrt{\log\mathcal{N}(\mathcal{I},\epsilon,\|\cdot\|)}d\epsilon\right)\end{align*}

\begin{align*}
		&\int_{\alpha}^{2e\sqrt{m}}\sqrt{\log\mathcal{N}(\mathcal{I},\epsilon,\|\cdot\|)}d\epsilon \\
		&\leq\int_{\alpha}^{2e\sqrt{m}}\sqrt{\log\mathcal{N}(\mathcal{B},\frac{\epsilon}{2},dist(\cdot,\cdot))}d\epsilon \\
		&\leq\int_{\alpha}^{2e\sqrt{m}}\sqrt{\log U}\leq2e\sqrt{m}\sqrt{\operatorname{log}U}
		\quad=2\sqrt{m\operatorname{log}U}
\end{align*}

\begin{align*}2|R|r^2\log\left(\frac{16M\sqrt r\max\{B_xB_1,\overline{R}B_2\}}{\alpha}
	\right)
\end{align*}

}
%\end{verbatim}

\end{document}
%%
%% End of file `sample-sigconf.tex'.
